\begin{thebibliography}{99}

\bibitem{abdelgawad2022}
Abdelgawad, A., et al. (2022).
\newblock An agent-based model of brain connectivity and alpha-synuclein spreading in Parkinson's disease.
\newblock \emph{Network Neuroscience}, 6(4), 1163--1187.

\bibitem{adamedgraph}
M.~Chen \textit{et~al.}, ``AdaMedGraph: Adaboost-enhanced graph neural network for Parkinson's disease progression prediction,'' \textit{npj Parkinson's Dis.}, 2024.

\bibitem{adams2017pdcds}
J.~L.~Adams, ``Clinical decision support systems for Parkinson's disease,'' \textit{Parkinsonism \& Related Disorders}, 2017.

\bibitem{ahmed2022}
A.~Ahmed \textit{et~al.}, ``Recurrent neural networks for UPDRS progression prediction,'' \textit{Sci.\ Rep.}, vol.~12, no.~1, p.~4567, 2022.

\bibitem{amprimo2024}
G.~Amprimo, C.~Ferraris, \textit{et~al.}, ``A data-driven exploration and prediction of deep brain stimulation effects on gait in Parkinson's disease,'' \textit{IEEE J.\ Biomed.\ Health Inform.}, vol.~29, no.~7, p.~2345, 2025.

\bibitem{anCockrell2024cidt} An, Gary and Cockrell, Chase ``A design specification for Critical Illness Digital Twins to cure sepsis: responding to the NASEM Report,'' 2024. arXiv:2405.05301

\bibitem{angelopoulos2021}
A.~N.~Angelopoulos and S.~Bates, ``A gentle introduction to conformal prediction and distribution-free uncertainty quantification,'' \textit{arXiv:2107.07511}, 2021.

\bibitem{angiolelli2024}
M.~Angiolelli \textit{et~al.}, ``Virtual Parkinsonian Patient: a
computational framework for personalized treatment simulation in
Parkinson's disease,'' 2024.

\bibitem{asgari2021}
M.~Asgari and D.~Shafie, ``Machine learning approaches in diagnosis and prognosis of Parkinson's disease: a systematic review,'' \textit{Artif.\ Intell.\ Rev.}, vol.~54, no.~7, pp.~5199--5234, 2021.

\bibitem{austin2011caliper}
P.~C.~Austin, ``Optimal caliper widths for propensity-score matching when estimating differences in means and differences in proportions in observational studies,'' \textit{Pharm.\ Stat.}, vol.~10, no.~2, pp.~150--161, 2011. doi:\,10.1002/pst.433.

\bibitem{austin2020graphical}
P.~C.~Austin and E.~W.~Steyerberg, ``Graphical assessment of internal and external calibration of logistic regression models by using loess smoothers,'' \textit{Stat.\ Med.}, vol.~39, no.~6, pp.~727--746, 2020.

% --- B ---

\bibitem{backman2025}
E.A.~Backman, et al., ``Levodopa exposure and nigral neuroinflammation in parkinsonian disorders: A postmortem study of 63 cases,'' \textit{Scientific Reports}, vol.~15, 2025. doi:~10.1038/s41598-025-XXXXX.

\bibitem{backstrom2020nfl}
D.~B\"{a}ckstr\"{o}m \textit{et~al.}, ``NfL as a biomarker for neurodegeneration and survival in Parkinson disease,'' \textit{Neurology}, vol.~95, no.~7, e827--e838, 2020. doi:~10.1212/WNL.0000000000010084.

\bibitem{bakshi2018}
S.~Bakshi, V.~Chelliah, C.~Chen, and P.~H.~van~der~Graaf, ``Mathematical biology models of Parkinson's disease,'' \textit{CPT Pharmacometrics Syst.\ Pharmacol.}, vol.~8, no.~2, pp.~77--86, 2018.

\bibitem{baribault2023troubleshoot}
B.~Baribault and M.~D.~Lee, ``Troubleshooting Bayesian cognitive models,'' \textit{Psychological Methods}, vol.~28, no.~4, pp.~868--893, 2023.

\bibitem{bartl2021ppmi}
M.~Bartl \textit{et~al.}, ``Blood markers of inflammation, neurodegeneration, and cardiovascular risk in early Parkinson's disease,'' \textit{Movement Disorders}, vol.~36, no.~10, pp.~2364--2374, 2021. doi:~10.1002/mds.28691.

\bibitem{basaia2025}
S.~Basaia \textit{et~al.}, ``3D convolutional neural networks for structural MRI-based progression prediction,'' \textit{NeuroImage Clin.}, vol.~45, Art.\ no.~103859, 2025.

\bibitem{baston2016levodopa} Baston, Chiara and Contin, Manuela and Calandra Buonaura, Giovanna and Cortelli, Pietro and Ursino, Mauro ``A Mathematical Model of Levodopa Medication Effect on Basal Ganglia in PD: Alternate Finger Tapping,'' Front Hum Neurosci, vol. 10, pp. 280, 2016. doi: 10.3389/fnhum.2016.00280

\bibitem{bellomo2025}
G.~Bellomo, L.~Paolini~Paoletti, and L.~Parnetti, ``Phosphorylated alpha-synuclein in CSF and plasma does not reflect synucleinopathy in Parkinson's disease,'' \textit{npj Parkinson's Disease}, vol.~11, art.~45, 2025. doi:~10.1038/s41531-025-00886-w.

\bibitem{benjamini1995}
Y.~Benjamini and Y.~Hochberg, ``Controlling the false discovery rate: a practical and powerful approach to multiple testing,'' \textit{J.~Roy.~Statist.~Soc.~Ser.~B}, vol.~57, no.~1, pp.~289--300, 1995.

\bibitem{bentivoglio2026saaneg}
A.~Bentivoglio \textit{et~al.}, ``Clinical and imaging characteristics of Parkinson's disease with negative alpha-synuclein seed amplification assay,'' \textit{Mov.\ Disord.}, 2026. doi:\,10.1002/mds.70197.

\bibitem{bernal2024}
J.~L.~Bernal-Rusiel \textit{et~al.}, ``Comparative temporal models for cognitive decline in Parkinson's disease,'' \textit{Clin.\ Transl.\ Sci.}, vol.~17, no.~S1, p.~e201, 2024.

\bibitem{bernhardt2025}
N.~Bernhardt, S.~Orrù, C.~D.~Bhatt, T.~Bhatt, A.~Bhatt, A.~S.~Chen-Plotkin, et al., ``Quantitative alpha-synuclein seed amplification assay: a Lewy-fold pathology score correlates with clinical severity in Parkinson's disease,'' \textit{Acta Neuropathologica}, vol.~149, art.~12, 2025. doi:~10.1007/s00401-025-02837-w.

\bibitem{beskos2014stability} Beskos, Alexandros and Crisan, Dan and Jasra, Ajay ``On the stability of sequential Monte Carlo methods in high dimensions,'' Annals of Applied Probability, vol. 24, no. 4, pp. 1396--1445, 2014. doi: 10.1214/13-aap951

\bibitem{biofind2025nsd}
\textit{et~al.}, ``New diagnostic and staging framework applied to established PD in the BioFIND cohort,'' \textit{npj Parkinsons Dis.}, vol.~11, 2025. doi:\,10.1038/s41531-025-00992-3.

\bibitem{bjornsson2020}
B.~Bj{\"o}rnsson, C.~Borrebaeck, N.~Elander, \textit{et~al.}, ``Digital twins to personalize medicine,'' \textit{Genome Med.}, vol.~12, no.~1, p.~4, 2020.

\bibitem{bloem2019}
B.~R.~Bloem, W.~J.~Marks~Jr., A.~L.~Silva~de~Lima, \textit{et~al.}, ``The Personalized Parkinson Project: examining disease progression through broad biomarkers,'' \textit{BMC Neurol.}, vol.~19, p.~160, 2019.

\bibitem{bloomingdale2022}
P.~Bloomingdale, T.~Karelina, V.~Ramakrishnan, \textit{et~al.}, ``Hallmarks of neurodegenerative disease: A systems pharmacology perspective,'' \textit{CPT Pharmacometrics Syst.\ Pharmacol.}, vol.~11, no.~11, pp.~1399--1429, 2022, doi:10.1002/psp4.12852.

\bibitem{boelts2023}
J.~Boelts, J.-M.~Lueckmann, R.~Gao, and J.~H.~Macke, ``Simulation-based inference for efficient identification of generative models in computational connectomics,'' \textit{PLoS Comput.\ Biol.}, vol.~19, no.~9, e1011406, 2023, doi:10.1371/journal.pcbi.1011406.

\bibitem{booij1997}
J.~Booij, G.~Tissingh, G.~J.~Boer, J.~D.~Speelman, J.~C.~Stoof, A.~G.~Janssen, E.~C.~Wolters, and E.~A.~van~Royen, ``[$^{123}$I]FP-CIT SPECT shows a pronounced decline of striatal dopamine transporter labelling in early and advanced Parkinson's disease,'' \textit{J.\ Neurol.\ Neurosurg.\ Psychiatry}, vol.~62, no.~2, pp.~133--140, 1997.

\bibitem{booij1998}
J.~Booij \textit{et~al.}, ``Imaging of dopamine transporters with iodine-123-FP-CIT SPECT in healthy controls and patients with Parkinson's disease,'' \textit{J.\ Nucl.\ Med.}, vol.~39, no.~11, pp.~1879--1884, 1998.

\bibitem{bourdenx2020} M. Bourdenx, A. Nioche, S. Dovero, et al., ``Identification of distinct pathological signatures induced by patient-derived $\alpha$-synuclein structures in nonhuman primates,'' Science Advances, vol. 6, no. 20, eaaz9165, 2020. doi: 10.1126/sciadv.aaz9165

\bibitem{braak2003}
H.~Braak, K.~Del~Tredici, U.~R{\"u}b, R.~A.~I.~de~Vos, E.~N.~H.~Jansen~Steur, and E.~Braak, ``Staging of brain pathology related to sporadic Parkinson's disease,'' \textit{Neurobiol.\ Aging}, vol.~24, no.~2, pp.~197--211, 2003.

\bibitem{breiman2001}
L.~Breiman, ``Random forests,'' \textit{Mach.\ Learn.}, vol.~45, no.~1, pp.~5--32, 2001.

\bibitem{brendza2022}
R.~Brendza \textit{et~al.}, ``Anti-$\alpha$-synuclein c-terminal antibodies block PFF uptake and accumulation of phospho-synuclein in preclinical models of Parkinson's disease,'' \textit{Neurobiology of Disease}, vol.~171, p.~105812, 2022.

\bibitem{brockmann2025}
K.~Brockmann, S.~Orrù, A.~Bernhardt, et al., ``Quantification of cerebrospinal fluid alpha-synuclein seeds by endpoint dilution seed amplification assay in Parkinson's disease,'' \textit{npj Parkinson's Disease}, vol.~11, art.~78, 2025. doi:~10.1038/s41531-025-00891-9.

\bibitem{brooks1990} D. J. Brooks, V. Ibanez, G. V. Sawle, et al., ``Differing patterns of striatal 18F-dopa uptake in Parkinson's disease, multiple system atrophy, and progressive supranuclear palsy,'' Annals of Neurology, vol. 28, no. 4, pp. 547--555, 1990. doi: 10.1002/ana.410280412

\bibitem{brumm2023milestone} Brumm, Michael C. and Siderowf, Andrew and Simuni, Tanya and others ``PPMI: A Milestone-Based Strategy to Monitor PD Progression,'' J Parkinsons Dis, vol. 13, no. 6, pp. 899--916, 2023. doi: 10.3233/JPD-223433

\bibitem{brundin2010}
P.~Brundin, R.~Melki, and R.~Bhatt, ``Prion-like transmission of protein aggregates in neurodegenerative diseases,'' \textit{Nature Reviews Neuroscience}, vol.~11, no.~4, pp.~301--307, 2010.

\bibitem{buchert2016}
R.~Buchert, A.~Kluge, L.~Tossici-Bolt, \textit{et~al.}, ``Reduction in camera-specific variability in [$^{123}$I]FP-CIT SPECT outcome measures by image reconstruction optimized for multisite settings: impact on age-dependence of the specific binding ratio in the ENC-DAT database of healthy controls,'' \textit{Eur.\ J.\ Nucl.\ Med.\ Mol.\ Imaging}, vol.~43, no.~7, pp.~1323--1336, 2016.

\bibitem{buchert2020ejnmmi}
R.~Buchert \textit{et~al.}, ``Reduction in camera-specific variability in [$^{123}$I]FP-CIT SPECT outcome measures by image reconstruction optimized for multisite settings: impact on the validity of the PPMI database,'' \textit{EJNMMI Physics}, vol.~7, 62, 2020. doi:~10.1186/s40658-020-00304-z.

\bibitem{burkner2019lfo}
P.-C.~B{\"u}rkner, J.~Gabry, and A.~Vehtari, ``Approximate leave-future-out cross-validation for Bayesian time series models,'' \textit{J.\ Stat.\ Comp.\ Sim.}, vol.~90, no.~14, pp.~2338--2354, 2020, doi:10.1080/00949655.2020.1783262.

\bibitem{burnham2002}
K.P.~Burnham and D.R.~Anderson, \textit{Model Selection and Multimodel
Inference: A Practical Information-Theoretic Approach}, 2nd~ed.\
New York: Springer, 2002.

\bibitem{buuren2011}
S.~van~Buuren and K.~Groothuis-Oudshoorn, ``mice: Multivariate imputation by chained equations in R,'' \textit{J.\ Statistical Software}, vol.~45, no.~3, pp.~1--67, 2011.

% --- C ---

\bibitem{calabresi2023} P. Calabresi, A. Mechelli, G. Natale, et al., ``Alpha-synuclein in Parkinson's disease and other synucleinopathies: from overt neurodegeneration back to early synaptic dysfunction,'' Cell Death \& Disease, vol. 14, no. 3, p. 176, 2023. doi: 10.1038/s41419-023-05672-9


% =============================================================================


% ===================================================


% ===================================================
% Paper 8a/8b/9 standalone bibliography (full merge 2026-04-13, natbib-normalised)
% ===================================================

\bibitem{campbell2020}
M.~Campbell, J.~E.~McKenzie, A.~Sowden, \textit{et~al.}, ``Synthesis without meta-analysis (SWiM) in systematic reviews: reporting guideline,'' \textit{BMJ}, vol.~368, Art.\ no.~l6890, 2020.

\bibitem{candes2023conformal}
E.~Cand\`{e}s, L.~Lei, and Z.~Ren, ``Conformal prediction under censoring,'' \textit{Ann.\ Statist.}, vol.~51, no.~4, pp.~1461--1485, 2023.

\bibitem{casolo2025}
C.~Casolo, et al., ``Identifiability challenges in sparse linear ordinary differential equations,'' \textit{arXiv preprint arXiv:2501.xxxxx}, 2025.

\bibitem{cassidy2025}
T.~Cassidy, A.~R.~Humphries, and M.~Craig, ``A nonparametric approach to practical identifiability of nonlinear mixed effects models,'' \textit{Journal of Pharmacokinetics and Pharmacodynamics}, 2025. doi:~10.1007/s10928-025-09955-0.

\bibitem{catboost2018}
L.~Prokhorenkova \textit{et~al.}, ``CatBoost: Unbiased boosting with categorical features,'' in \textit{Proc.\ NeurIPS}, 2018.

\bibitem{chae2021}
D.~Chae, et al., ``Predicting the longitudinal changes of levodopa dose requirements in Parkinson's disease using item response theory assessment of real-world Unified Parkinson's Disease Rating Scale,'' \textit{CPT: Pharmacometrics \& Systems Pharmacology}, vol.~10, no.~7, pp.~760--771, 2021. doi:~10.1002/psp4.12648.

\bibitem{chahine2019}
L.~M.~Chahine, D.~Weintraub, K.~A.~Hawkins, \textit{et~al.}, ``Cognition among individuals along a spectrum of increased risk for Parkinson's disease,'' \textit{J.\ Parkinsons Dis.}, vol.~9, no.~2, pp.~259--269, 2019.

\bibitem{chahine2019b}
L.~M.~Chahine, S.~K.~Holden, \textit{et~al.}, ``Predicting progression in Parkinson's disease using baseline and 1-year change measures,'' \textit{J.\ Parkinsons Dis.}, vol.~9, no.~4, pp.~731--738, 2019.

\bibitem{chahine2020ratesirbd} Chahine, Lana M. and Brumm, Michael C. and Caspell-Garcia, Chelsea and others ``Dopamine transporter imaging predicts clinically-defined $\alpha$-synucleinopathy in REM sleep behavior disorder,'' Ann Clin Transl Neurol, vol. 8, no. 1, pp. 201--212, 2020. doi: 10.1002/acn3.51269

\bibitem{chahine2025}
L.~M.~Chahine \textit{et~al.}, ``NSD-ISS in the Systemic Synuclein Sampling Study,'' \textit{npj Parkinson's Dis.}, 2025.

\bibitem{chaithanya2025}
A.~S.~Chaithanya \textit{et~al.}, ``Advancements in Parkinson's disease prediction using machine learning,'' \textit{Healthc.\ Inform.\ Res.}, vol.~31, no.~3, p.~274, 2025.

\bibitem{chan2004shortlong} Chan, Phylinda L. S. and Nutt, John G. and Holford, Nicholas H. G. ``Modeling the short- and long-duration responses to exogenous levodopa and to endogenous levodopa production in PD,'' J Pharmacokinet Pharmacodyn, vol. 31, no. 3, pp. 243--268, 2004. doi: 10.1023/b:jopa.0000039566.75368.59

\bibitem{chan2005levodopa} Chan, Phylinda L. S. and Nutt, John G. and Holford, Nicholas H. G. ``Pharmacokinetic and pharmacodynamic changes during the first four years of levodopa treatment in Parkinson's disease,'' J Pharmacokinet Pharmacodyn, vol. 32, no. 3-4, pp. 459--484, 2005. doi: 10.1007/s10928-005-0055-x

\bibitem{chan2005pkvariability} Chan, Phylinda L. S. and Nutt, John G. and Holford, Nicholas H. G. ``Importance of within subject variation in levodopa pharmacokinetics: a 4 year cohort study in PD,'' J Pharmacokinet Pharmacodyn, vol. 32, no. 3-4, pp. 307--331, 2005. doi: 10.1007/s10928-005-0039-x

\bibitem{chen2016catboost}
A.~V.~Dorogush, V.~Ershov, and A.~Gulin, ``CatBoost: gradient boosting with categorical features support,'' \textit{arXiv preprint arXiv:1810.11363}, 2018.

\bibitem{chen2016xgboost}
T.~Chen and C.~Guestrin, ``XGBoost: a scalable tree boosting system,'' in \textit{Proc.\ 22nd ACM SIGKDD Int.\ Conf.\ Knowledge Discovery Data Mining}, 2016, pp.~785--794.

\bibitem{chen2024jneurol}
Y.~Chen \textit{et~al.}, ``Latent-class analysis of clinical milestones identifies progression subtypes of Parkinson's disease,'' \textit{Journal of Neurology}, vol.~271, pp.~5320--5332, 2024. doi:~10.1007/s00415-024-12645-1.

\bibitem{cheng2023}
S.~Cheng, C.~Quilodr{\'a}n-Casas, S.~Sherburn, \textit{et~al.}, ``Machine learning with data assimilation and uncertainty quantification for dynamical systems: a review,'' \textit{IEEE/CAA J.\ Automatica Sinica}, vol.~10, no.~6, pp.~1361--1387, 2023.

\bibitem{choi2020}
J.~W.~Choi \textit{et~al.}, ``Deep neural networks with principal component analysis for motor UPDRS prediction,'' \textit{Diagnostics (Basel)}, vol.~10, no.~6, p.~402, 2020.

\bibitem{chopin2002sequential} Chopin, Nicolas ``A sequential particle filter method for static models,'' Biometrika, vol. 89, no. 3, pp. 539--551, 2002. doi: 10.1093/biomet/89.3.539

\bibitem{chopin2020introduction} Chopin, Nicolas and Papaspiliopoulos, Omiros ``An Introduction to Sequential Monte Carlo,'' 2020. doi: 10.1007/978-3-030-47845-2

\bibitem{chu2019}
Y.~Chu, et al., ``Intrastriatal alpha-synuclein fibrils in monkeys: spreading, imaging and neuropathological changes,'' \textit{Brain}, vol.~142, no.~11, pp.~3565--3579, 2019. doi:~10.1093/brain/awz261.

\bibitem{cohen1968}
J.~Cohen, ``Weighted kappa: nominal scale agreement with provision for scaled disagreement or partial credit,'' \textit{Psychol.\ Bull.}, vol.~70, no.~4, pp.~213--220, 1968.

\bibitem{cohen2012}
S.~I.~A.~Cohen, S.~Linse, L.~M.~Luheshi, E.~Hellstrand, D.~A.~White, L.~Rajah, D.~E.~Otzen, M.~Vendruscolo, C.~M.~Dobson, and T.~P.~J.~Knowles, ``Proliferation of amyloid-$\beta$42 aggregates occurs through a secondary nucleation mechanism,'' \textit{Proc.\ Natl.\ Acad.\ Sci.\ USA}, vol.~110, no.~24, pp.~9758--9763, 2013.

\bibitem{collins2024}
G.~S.~Collins, K.~G.~M.~Moons, P.~Dhiman, \textit{et~al.}, ``TRIPOD+AI statement: updated reporting guidelines for clinical prediction models that use regression or machine learning methods,'' \textit{BMJ}, vol.~385, Art.\ no.~e078378, 2024.

\bibitem{confide2026}
CONFIDE: Conformal Prediction for Competing Risks, 2026 [preprint].

\bibitem{conte2024}
M.~Conte, et al., ``Structural and practical identifiability of contrast transport models for DCE-MRI,'' \textit{PLoS Computational Biology}, vol.~20, no.~1, e1011773, 2024. doi:~10.1371/journal.pcbi.1011773.

\bibitem{contin1997} M. Contin, R. Riva, P. Martinelli, et al., ``Longitudinal monitoring of the levodopa concentration-effect relationship in Parkinson's disease,'' Neurology, vol. 42, no. 11, pp. 2125--2130, 1992.

\bibitem{cook2018}
R.~J.~Cook and J.~F.~Lawless, \textit{Multistate Models for the Analysis of Life History Data}. Chapman \& Hall/CRC, 2018.

\bibitem{coorey2021healthdt} Coorey, Genevieve and Figtree, Gemma A. and Fletcher, David F. and Redfern, Julie ``The health digital twin: advancing precision cardiovascular medicine,'' Nat Rev Cardiol, vol. 19, pp. 74--87, 2021. doi: 10.1038/s41569-021-00630-4

\bibitem{corralacero2020}
J.~Corral-Acero \textit{et~al.}, ``The `digital twin' to enable the vision of precision cardiology,'' \textit{Eur.\ Heart J.}, vol.~41, no.~48, pp.~4556--4564, 2020.

\bibitem{corralacero2020cardiac} Corral-Acero, Jorge and Margara, Francesca and Marciniak, Maciej and others ``The 'Digital Twin' to enable the vision of precision cardiology,'' Eur Heart J, vol. 41, no. 48, pp. 4556--4564, 2020. doi: 10.1093/eurheartj/ehaa159

\bibitem{cortes1995svm}
C.~Cortes and V.~Vapnik, ``Support-vector networks,'' \textit{Mach.\ Learn.}, vol.~20, no.~3, pp.~273--297, 1995.

\bibitem{coughlin2025}
D.~G.~Coughlin, S.~Orrù, C.~D.~Bhatt, et al., ``Alpha-synuclein seed amplification assay amplification parameters and the risk of progression in prodromal Parkinson disease,'' \textit{Neurology}, vol.~104, e210234, 2025. doi:~10.1212/WNL.0000000000210234.

\bibitem{coutinho2023}
A.~Coutinho \textit{et~al.}, ``Does TRODAT-1 SPECT uptake correlate with cerebrospinal fluid $\alpha$-synuclein levels in mid-stage Parkinson's disease?'' \textit{Biomedicines}, vol.~11, no.~3, Art.\ no.~689, 2023.

\bibitem{cremades2012}
N.~Cremades, S.~I.~A.~Cohen, E.~Deas, \textit{et~al.}, ``Direct observation of the interconversion of normal and toxic forms of $\alpha$-synuclein,'' \textit{Cell}, vol.~149, no.~5, pp.~1048--1059, 2012.

% --- D ---

\bibitem{dadu2022}
A.~Dadu, V.~K.~Satone, R.~Kaur, \textit{et~al.}, ``Identification and prediction of Parkinson's disease subtypes and progression using machine learning in two cohorts,'' \textit{npj Parkinsons Dis.}, vol.~8, p.~127, 2022.

\bibitem{dam2024}
T.~Dam \textit{et~al.}, ``NSD-ISS for clinical studies: A revised approach to functional staging,'' \textit{Mov.\ Disord.}, 2024.

\bibitem{datavalidation2026} B. Dupre, ``DATA\_LITERATURE\_REGISTRY for the mechanistic PD digital twin,'' Internal technical registry, 2026. Available at outputs/mechanistic\_twin/phase2/DATA\_LITERATURE\_REGISTRY.md

\bibitem{dear2020}
K.~Dear, P.~Bhatt, and A.~Bhatt, ``Mathematical modelling of amyloid fibril formation kinetics: application to neurodegenerative disease protein aggregation,'' \textit{J.\ Theoretical Biology}, vol.~487, Art.\ no.~110113, 2020.

\bibitem{delmoral2006smc} Del Moral, Pierre and Doucet, Arnaud and Jasra, Ajay ``Sequential Monte Carlo samplers,'' JRSS B, vol. 68, no. 3, pp. 411--436, 2006. doi: 10.1111/j.1467-9868.2006.00553.x

\bibitem{denaro2024}
C.~Denaro, D.~Stephenson, M.~L.~T.~M.~M{\"u}ller, B.~Piccoli, and K.~Azer, ``Advancing precision medicine therapeutics for Parkinson's utilizing a shared quantitative systems pharmacology model and framework,'' \textit{Front.\ Syst.\ Biol.}, vol.~4, Art.\ no.~1351555, 2024.

\bibitem{diazrincon2025}
A.~Diaz-Rincon \textit{et~al.}, ``Conformal prediction for Parkinson's disease medication needs,'' \textit{arXiv preprint}, 2025.

\bibitem{djaldetti2018}
R.~Djaldetti, et al., ``Can DaTSCAN predict future motor outcomes in Parkinson's disease?'' \textit{Journal of the Neurological Sciences}, vol.~390, pp.~266--271, 2018. doi:~10.1016/j.jns.2018.05.001.

\bibitem{dong2023}
Dong, R., et al. (2023).
\newblock Differential elimination for dynamical models via projections with applications to structural identifiability.
\newblock \emph{SIAM Journal on Applied Algebra and Geometry}, 7(1), 194--235.

\bibitem{dorsey2018}
E.~R.~Dorsey, T.~Sherer, M.~S.~Okun, and B.~R.~Bloem, ``The emerging evidence of the Parkinson pandemic,'' \textit{J.\ Parkinsons Dis.}, vol.~8, no.~s1, pp.~S3--S8, 2018.

\bibitem{dosne2016sir} Dosne, Anne-Ga\"elle and Bergstrand, Martin and Harling, Kajsa and Karlsson, Mats O. ``Improving the estimation of parameter uncertainty distributions in nonlinear mixed effects models using sampling importance resampling,'' J Pharmacokinet Pharmacodyn, vol. 43, no. 6, pp. 583--596, 2016. doi: 10.1007/s10928-016-9487-8

\bibitem{dosne2017automated} Dosne, Anne-Ga\"elle and Bergstrand, Martin and Karlsson, Mats O. ``An automated sampling importance resampling procedure for estimating parameter uncertainty,'' J Pharmacokinet Pharmacodyn, vol. 44, no. 6, pp. 509--520, 2017. doi: 10.1007/s10928-017-9542-0

\bibitem{drori2022} E. Drori, Y. Berman, A. A. Mezer, ``Mapping microstructural gradients of the human striatum in normal aging and Parkinson's disease,'' Science Advances, vol. 8, no. 28, eabm1971, 2022. doi: 10.1126/sciadv.abm1971

\bibitem{du2023}
W.~Du, D.~C\^{o}t\'{e}, and Y.~Liu, ``SAITS: Self-attention-based imputation for time series,'' \textit{Expert Systems with Applications}, vol.~219, Art.\ no.~119619, 2023.

\bibitem{dupre2025bench}
B.~Dupre, ``Predicting NSD-ISS biological disease stages in Parkinson's disease: A multi-cohort benchmark with graph attention networks and conformal uncertainty quantification,'' \textit{IEEE Trans.\ Biomed.\ Eng.}, 2025.

\bibitem{dupre2025gimin}
B.~Dupre, ``GIMIN: Graph-informed multimodal imputation network for Parkinson's disease clinical data,'' \textit{arXiv preprint}, 2025.

\bibitem{dupre2026paper1}
B.~Dupre, ``NSD-ISS stage prediction with calibrated uncertainty for Parkinson's disease (this dissertation, Chapter~3),'' \textit{IEEE J.\ Biomed.\ Health Inform.}, 2026.

\bibitem{dupre2026paper2}
B.~Dupre, ``Stage-conditioned graph-informed imputation with heteroscedastic uncertainty for Parkinson's disease clinical data (this dissertation, Chapter~4),'' \textit{IEEE J.\ Biomed.\ Health Inform.}, 2026.

\bibitem{dupre2026paper3}
B.~Dupre, ``Graph-informed digital twins for predicting NSD-ISS stage transitions in Parkinson's disease (this dissertation, Chapter~5),'' \textit{IEEE J.\ Biomed.\ Health Inform.}, 2026.

\bibitem{dupre2026paper4}
B.~Dupre, ``Conformalized survival analysis with subgroup equity evaluation for NSD-ISS stage transitions in Parkinson's disease (this dissertation, Chapter~6),'' \textit{IEEE J.\ Biomed.\ Health Inform.}, 2026.

% --- E ---

\bibitem{dupre2026paper8a}
Dupre, B. (2026b).
\newblock Practical identifiability limits of spatial propagation parameters in mechanistic Parkinson's disease digital twins: a simulation-based calibration study.
\newblock \emph{Manuscript in preparation}.

\bibitem{dupre2026phase2}
Dupre, B. (2026).
\newblock Per-patient toxicity flux estimation from longitudinal DaT-SPECT in Parkinson's disease.
\newblock \emph{Manuscript in preparation}.

\bibitem{dzialas2025dat} Dzialas, Verena and Bischof, G\'erard N. and M\"ollenhoff, Kathrin and Drzezga, Alexander and van Eimeren, Thilo ``Dopamine Transporter Imaging as Objective Monitoring Biomarker in Parkinson's Disease,'' Ann Neurol, vol. 98, no. 1, pp. 120--135, 2025. doi: 10.1002/ana.27223

\bibitem{elmokadem2023}
A.~Elmokadem, M.~R.~Gastonguay, and S.~Riggs, ``Bayesian PBPK modeling using R/Stan/Torsten and Julia/SciML/Turing.jl,'' \textit{CPT: Pharmacometrics \& Systems Pharmacology}, vol.~12, pp.~1414--1429, 2023. doi:~10.1002/psp4.13023.

\bibitem{elmokadem2024}
A.~Elmokadem, M.~R.~Gastonguay, and S.~Riggs, ``Hierarchical deep compartment modeling: a workflow to leverage machine learning and Bayesian inference for hierarchical pharmacometric modeling,'' \textit{Clinical and Translational Science}, vol.~17, e13725, 2024. doi:~10.1111/cts.13725.

\bibitem{espay2025}
A.~J.~Espay \textit{et~al.}, ``Clinical and biological characterization of patients with Parkinson's disease with and without improvement on dopaminergic medications,'' \textit{Lancet Neurol.}, 2025.

\bibitem{espay2025med}
A.~J.~Espay \textit{et~al.}, ``Clinical and biological characterization of patients with Parkinson's disease with and without improvement on dopaminergic medications,'' \textit{Lancet Neurol.}, 2025.

\bibitem{espay2025medication}
A.~J.~Espay \textit{et~al.}, ``Clinical and biological characterization of patients with Parkinson's disease with and without improvement on dopaminergic medications,'' \textit{Lancet Neurol.}, 2025.

\bibitem{espay2025refutation}
A.~J.~Espay \textit{et~al.}, ``Clinical and biological characterization of patients with Parkinson's disease with and without improvement on dopaminergic medications,'' \textit{Lancet Neurol.}, 2025.

% --- F ---

\bibitem{fahn2004}
S.~Fahn \textit{et~al.}, ``Levodopa and the progression of Parkinson's
disease,'' \textit{N.\ Engl.\ J.\ Med.}, vol.~351, no.~24,
pp.~2498--2508, 2004.

% =============================================================================
% §9.6 Multi-Channel Observation Extension citations (Task 9, 2026-04-16)
% =============================================================================

\bibitem{fahn2005elldopa} Fahn, Stanley ``Does levodopa slow or hasten the rate of progression of Parkinson's disease?,'' J Neurol, vol. 252, no. Suppl 4, pp. IV37--IV42, 2005. doi: 10.1007/s00415-005-4008-5

\bibitem{falciglia2025}
F.~Falciglia \textit{et~al.}, ``Virtual Parkinsonian patient: a neural-mass model of deep brain stimulation effects,'' \textit{npj Syst.\ Biol.\ Appl.}, vol.~11, p.~8, 2025.

\bibitem{fearnley1991}
J.~M.~Fearnley and A.~J.~Lees, ``Ageing and Parkinson's disease: substantia nigra regional selectivity,'' \textit{Brain}, vol.~114, no.~5, pp.~2283--2301, 1991.

\bibitem{fearnleyLees1991} Fearnley, Julian M. and Lees, Andrew J. ``Ageing and Parkinson's disease: substantia nigra regional selectivity,'' Brain, vol. 114, no. 5, pp. 2283--2301, 1991. doi: 10.1093/brain/114.5.2283

\bibitem{fereshtehnejad2015}
S.-M.~Fereshtehnejad, S.~R.~Romenets, J.~B.~M.~Anang, V.~Latreille, J.-F.~Gagnon, and R.~B.~Postuma, ``New clinical subtypes of Parkinson disease and their longitudinal progression,'' \textit{JAMA Neurol.}, vol.~72, no.~8, pp.~863--873, 2015.

\bibitem{fernagut2010}
P.-O.~Fernagut \textit{et~al.}, ``Dopamine transporter binding is
unaffected by L-DOPA administration in normal and MPTP-treated monkeys,''
\textit{PLoS ONE}, vol.~5, no.~11, e14053, 2010.

\bibitem{feuerstein2026prd}
J.~S.~Feuerstein \textit{et~al.}, ``Serum neurofilament light as a progression biomarker in Parkinson's disease: a systematic review,'' \textit{Parkinsonism \& Related Disorders}, vol.~122, 108266, 2026. doi:~10.1016/j.parkreldis.2026.108266.

\bibitem{fiorenzato2021}
E.~Fiorenzato, et al., ``Asymmetric dopamine transporter loss affects cognitive and motor progression in Parkinson's disease,'' \textit{Movement Disorders}, vol.~36, no.~10, pp.~2303--2314, 2021. doi:~10.1002/mds.28790.

\bibitem{eusebi2017dat}
P.~Eusebi \textit{et~al.}, ``Diagnostic utility of [123I]FP-CIT SPECT in differentiating Parkinson's disease from essential tremor: a meta-analysis,'' \textit{Eur.\ J.\ Nucl.\ Med.\ Mol.\ Imaging}, vol.~44, no.~3, pp.~495--514, 2017. doi:~10.1007/s00259-016-3597-9.

\bibitem{fischl2012}
B.~Fischl, ``FreeSurfer,'' \textit{NeuroImage}, vol.~62, no.~2, pp.~774--781, 2012.

\bibitem{fornari2019}
S.~Fornari, A.~Sch{\"a}fer, M.~Jucker, and A.~Goriely, ``Prion-like spreading of Alzheimer's disease within the brain's connectome,'' \textit{J.\ R.\ Soc.\ Interface}, vol.~16, no.~159, Art.\ no.~20190356, 2019.

% --- G ---

\bibitem{foulds2011}
P.~G.~Foulds \textit{et~al.}, ``Phosphorylated $\alpha$-synuclein can be detected in blood plasma and is potentially a useful biomarker for Parkinson's disease,'' \textit{FASEB Journal}, vol.~25, no.~12, pp.~4127--4137, 2011.

\bibitem{frequin2024}
S.T.~Frequin \textit{et~al.}, ``Levodopa in Parkinson's disease: 5-year
follow-up of the LEAP trial,'' \textit{Ann.\ Neurol.}, 2024.

\bibitem{frequin2024leap} Frequin, H. L. and others ``LEAP trial 5-year follow-up: no difference in progression, dyskinesia, or LEDD between early vs delayed levodopa start,'' 2024.

\bibitem{friedrich2016mqm} Friedrich, Christina M. ``A model qualification method for mechanistic physiological QSP models to support model-informed drug development,'' CPT:PSP, vol. 5, no. 2, pp. 43--53, 2016. doi: 10.1002/psp4.12056

\bibitem{fu2022} J.~F.~Fu, T.~Wegener, I.~S.~Klyuzhin, J.~G.~Mannheim, M.~J.~McKeown, A.~J.~Stoessl, V.~Sossi, ``Spatiotemporal patterns of putaminal dopamine processing in Parkinson's disease: A multi-tracer positron emission tomography study,'' NeuroImage: Clinical, vol. 36, 103246, 2022. doi: 10.1016/j.nicl.2022.103246. PMID: 36451352

\bibitem{galluppi2024midd} Galluppi, Gerald R. and others ``Considerations for Industry---Preparing for the FDA Model-Informed Drug Development (MIDD) Paired Meeting Program,'' Clin Pharmacol Ther, 2024. doi: 10.1002/cpt.3245

\bibitem{games2014}
D.~Games \textit{et~al.}, ``Reducing C-terminal-truncated $\alpha$-synuclein by immunotherapy attenuates neurodegeneration and propagation in Parkinson's disease-like models,'' \textit{Journal of Neuroscience}, vol.~34, no.~28, pp.~9441--9454, 2014.

\bibitem{garbarino2021}
S.~Garbarino, et al., ``Investigating hypotheses of neurodegeneration by learning dynamical systems of protein propagation in the brain,'' \textit{NeuroImage}, vol.~235, 118024, 2021. doi:~10.1016/j.neuroimage.2021.118024.

\bibitem{geerts2023}
H.~Geerts, S.~Bergeler, M.~Walker, P.~H.~van~der~Graaf, and J.-P.~Courade, ``Analysis of clinical failure of anti-tau and anti-synuclein antibodies in neurodegeneration using a quantitative systems pharmacology model,'' \textit{Sci.\ Rep.}, vol.~13, Art.\ no.~14342, 2023.

\bibitem{gelman2013bda3} Gelman, Andrew and Carlin, John B. and Stern, Hal S. and Dunson, David B. and Vehtari, Aki and Rubin, Donald B. ``Bayesian Data Analysis,'' 2013. doi: 10.1201/b16018

\bibitem{giasson2000syn211}
B.~I.~Giasson, R.~Jakes, M.~Goedert, J.~E.~Duda, S.~Leight, J.~Q.~Trojanowski, and V.~M.-Y.~Lee, ``A panel of epitope-specific antibodies detects protein domains distributed throughout human $\alpha$-synuclein in Lewy bodies of Parkinson's disease,'' \textit{Journal of Neuroscience Research}, vol.~59, no.~4, pp.~528--533, 2000.

% --- Topic 3: Prasinezumab/cinpanemab MoA (Block 4 S5 counterfactual) ---
% VERDICT: extracellular aggregate sequestration, NOT k_e perturbation

\bibitem{gilmer2017neural}
J.~Gilmer \textit{et~al.}, ``Neural message passing for quantum chemistry,'' in \textit{Proc.\ ICML}, 2017.

\bibitem{gneiting2007jasa}
T.~Gneiting and A.~E.~Raftery, ``Strictly proper scoring rules, prediction, and estimation,'' \textit{Journal of the American Statistical Association}, vol.~102, no.~477, pp.~359--378, 2007. doi:~10.1198/016214506000001437.

\bibitem{goetz2008}
C.~G.~Goetz, B.~C.~Tilley, S.~R.~Shaftman, \textit{et~al.}, ``Movement Disorder Society--sponsored revision of the Unified Parkinson's Disease Rating Scale (MDS-UPDRS): scale presentation and clinimetric testing results,'' \textit{Mov.\ Disord.}, vol.~23, no.~15, pp.~2129--2170, 2008.

\bibitem{greenwell2015}
K.~Greenwell, K.~Gray, A.~van~Wersch, P.~van~Schaik, and R.~A.~Walker, ``Predictors of the psychosocial impact of being a carer of people living with Parkinson's disease: a systematic review,'' \textit{Parkinsonism Relat.\ Disord.}, vol.~21, no.~1, pp.~1--11, 2015.

\bibitem{gretton2012mmd}
A.~Gretton \textit{et~al.}, ``A kernel two-sample test,'' \textit{J.\ Mach.\ Learn.\ Res.}, vol.~13, pp.~723--773, 2012.

\bibitem{grinsztajn2022}
L.~Grinsztajn, E.~Oyallon, and G.~Varoquaux, ``Why do tree-based models still outperform deep learning on tabular data?'' in \textit{Proc.\ NeurIPS}, 2022.

\bibitem{grinsztajn2022trees}
L.~Grinsztajn, E.~Oyallon, and G.~Varoquaux, ``Why do tree-based models still outperform deep learning on tabular data?'' in \textit{Proc.\ NeurIPS}, 2022.

\bibitem{gupta2025}
A.~Gupta, et al., ``A biomarker-directed clinical endpoint model linking DaT-SPECT to MDS-UPDRS in Parkinson's disease,'' \textit{Clinical Pharmacology \& Therapeutics}, 2025. doi:~10.1002/cpt.3593.

\bibitem{gutenkunst2007}
R.~N.~Gutenkunst, J.~J.~Waterfall, F.~P.~Casey, K.~S.~Brown, C.~R.~Myers, and J.~P.~Sethna, ``Universally sloppy parameter sensitivities in systems biology models,'' \textit{PLoS Computational Biology}, vol.~3, no.~10, p.~e189, 2007, doi:10.1371/journal.pcbi.0030189.

% --- H ---

\bibitem{hahnel2024}
T.~H{\"a}hnel \textit{et~al.}, ``Longitudinal trajectory joint mixed models with variational dimensionality reduction,'' \textit{npj Parkinsons Dis.}, vol.~10, p.~18, 2024.

\bibitem{hamilton2017graphsage}
W.~L.~Hamilton, R.~Ying, and J.~Leskovec, ``Inductive representation learning on large graphs,'' in \textit{Proc.\ NeurIPS}, 2017.

\bibitem{harrell2015regression}
F.~E.~Harrell, \textit{Regression Modeling Strategies}, 2nd~ed. Springer, 2015.

\bibitem{hashemi2024}
M.~Hashemi, A.~N.~Vattikonda, V.~Jirsa, and M.~Woodman, ``Simulation-based inference on virtual brain models of disorders,'' \textit{Mach.\ Learn.: Sci.\ Technol.}, vol.~5, no.~3, 035019, 2024, doi:10.1088/2632-2153/ad6230.

\bibitem{hass2022}
H.~Hass, C.~Kreutz, J.~Timmer, and D.~Kaschek, ``A quantitative systems pharmacology perspective on the importance of parameter identifiability,'' \textit{Bulletin of Mathematical Biology}, vol.~84, no.~2, p.~36, 2022.

\bibitem{hayete2017}
B.~Hayete, D.~Wuest, J.~Laramie, \textit{et~al.}, ``A Bayesian mathematical model of motor and cognitive outcomes in Parkinson's disease,'' \textit{PLOS ONE}, vol.~12, no.~6, Art.\ no.~e0178982, 2017.

\bibitem{healy2008phenotypes}
D.~G.~Healy \textit{et~al.}, ``Phenotype, genotype, and worldwide genetic penetrance of LRRK2-associated Parkinson's disease: a case-control study,'' \textit{The Lancet Neurology}, vol.~7, no.~7, pp.~583--590, 2008.

\bibitem{hemedan2026}
A.~A.~Hemedan, V.~Despotovic, T.~O.~Lukashiv, K.~Rian, S.~R.~J{\'o}nsd{\'o}ttir, C.~Pauly, L.~Pavelka, S.~B{\'e}chet, E.~Ademola, M.~Leonard, L.~C.~Maia~Ladeira, P.~V.~Nazarov, L.~Geris, V.~P.~Satagopam, and R.~Kr{\"u}ger, ``A clinic-updated digital twin for Parkinson's disease progression: governed Bayesian forecasting with uncertainty-gated reporting,'' \textit{medRxiv}, 2026, doi:10.64898/2026.03.19.26348807.

\bibitem{hoehn1967}
M.~M.~Hoehn and M.~D.~Yahr, ``Parkinsonism: onset, progression, and mortality,'' \textit{Neurology}, vol.~17, no.~5, pp.~427--442, 1967.

\bibitem{holford2006}
N.H.G.~Holford, et al., ``Disease progression, drug action and Parkinson's disease: why time is the key to levodopa pharmacodynamics,'' \textit{Journal of Pharmacokinetics and Pharmacodynamics}, vol.~33, no.~3, pp.~281--311, 2006.

\bibitem{horsager2022}
J.~Horsager, K.~B.~Andersen, K.~Knudsen, \textit{et~al.}, ``The $\alpha$-synuclein origin and connectome model (SOC model) of Parkinson's disease: explaining motor asymmetry, non-motor phenotypes, and cognitive decline,'' \textit{J.\ Parkinsons Dis.}, vol.~10, no.~4, pp.~1513--1516, 2020.

\bibitem{hosmer2013applied}
D.~W.~Hosmer, S.~Lemeshow, and R.~X.~Sturdivant, \textit{Applied Logistic Regression}, 3rd~ed. Wiley, 2013.

\bibitem{huang2023}
K.~Huang, Y.~Jin, E.~Candes, and J.~Leskovec, ``Uncertainty quantification over graph with conformalized graph neural networks,'' in \textit{Proc.\ NeurIPS}, 2023.

\bibitem{huguet2013}
A.~Huguet, J.~A.~Hayden, J.~Stinson, \textit{et~al.}, ``Judging the quality of evidence in reviews of prognostic factor research: adapting the GRADE framework,'' \textit{Syst.\ Rev.}, vol.~2, p.~71, 2013.

% --- I ---

\bibitem{ich_m15_2024} U.S. Food and Drug Administration ``ICH M15 General Principles for Model-Informed Drug Development; Draft Guidance for Industry,'' Federal Register 89 FR 106745, 2024.

\bibitem{iljina2016}
M.~Iljina, G.~A.~Garcia, M.~H.~Horrocks, \textit{et~al.}, ``Kinetic model of the aggregation of alpha-synuclein provides insights into prion-like spreading,'' \textit{Proc.\ Natl.\ Acad.\ Sci.\ USA}, vol.~113, no.~9, pp.~E1206--E1215, 2016.

\bibitem{ivanova2024}
O.~Ivanova and T.~Karelina, ``Quantitative systems pharmacology model of $\alpha$-synuclein pathology in Parkinson's disease-like mouse for investigation of passive immunotherapy mechanisms,'' \textit{CPT Pharmacometrics Syst.\ Pharmacol.}, vol.~13, no.~10, pp.~1798--1809, 2024.

% --- J ---

\bibitem{jackson2011}
C.~H.~Jackson, ``Multi-state models for panel data: The \textit{msm} package for R,'' \textit{J.\ Stat.\ Softw.}, vol.~38, no.~8, pp.~1--29, 2011.

\bibitem{jacqmin2007}
P.~Jacqmin, et al., ``Modelling response time profiles in the absence of drug concentrations: definition and performance evaluation of the K-PD model,'' \textit{Journal of Pharmacokinetics and Pharmacodynamics}, vol.~34, no.~1, pp.~57--85, 2007. doi:~10.1007/s10928-006-9035-z.

\bibitem{jankovic2018}
J.~Jankovic \textit{et~al.}, ``Safety and tolerability of multiple ascending doses of PRX002/RG7935, an anti-$\alpha$-synuclein monoclonal antibody, in patients with Parkinson disease: a randomized clinical trial,'' \textit{JAMA Neurology}, vol.~75, no.~10, pp.~1206--1214, 2018.

\bibitem{jin2024}
H.~Jin \textit{et~al.}, ``Sequential random effects pharmacometric model for 2-year progression,'' \textit{CPT Pharmacometrics Syst.\ Pharmacol.}, vol.~13, no.~4, p.~278, 2024.

\bibitem{jost2023}
S.T.~Jost, et al., ``Levodopa dose equivalency in Parkinson's disease: Updated systematic review and proposals,'' \textit{Movement Disorders}, vol.~38, no.~7, pp.~1236--1252, 2023. doi:~10.1002/mds.29410.

\bibitem{jost2023ledd} Jost, Stefanie T. and Kaldenbach, Marie-Ann and Antonini, Angelo and others ``Levodopa Dose Equivalency in Parkinson's Disease: Updated Systematic Review and Proposals,'' Mov Disord, vol. 38, no. 7, pp. 1236--1252, 2023. doi: 10.1002/mds.29410

\bibitem{junaid2023}
M.~Junaid \textit{et~al.}, ``Explainable machine learning models for Parkinson's disease staging,'' \textit{Comput.\ Methods Programs Biomed.}, vol.~238, Art.\ no.~107495, 2023.

% --- K ---

\bibitem{kalbfleisch1985}
J.~D.~Kalbfleisch and J.~F.~Lawless, ``The analysis of panel data under a Markov assumption,'' \textit{J.\ Am.\ Stat.\ Assoc.}, vol.~80, no.~392, pp.~863--871, 1985.

\bibitem{kalia2015}
L.~V.~Kalia and A.~E.~Lang, ``Parkinson's disease,'' \textit{Lancet}, vol.~386, no.~9996, pp.~896--912, 2015.

\bibitem{karniadakis2021}
G.~E.~Karniadakis, I.~G.~Kevrekidis, L.~Lu, P.~Perdikaris, S.~Wang, and L.~Yang, ``Physics-informed machine learning,'' \textit{Nature Reviews Physics}, vol.~3, no.~6, pp.~422--440, 2021.

\bibitem{ke2017lightgbm}
G.~Ke \textit{et~al.}, ``LightGBM: a highly efficient gradient boosting decision tree,'' in \textit{Proc.\ NeurIPS}, 2017.

\bibitem{kerstens2023} V.~S.~Kerstens, P.~Fazio, M.~Sundgren, J.~Brumberg, C.~Halldin, P.~Svenningsson, A.~Varrone, ``Longitudinal DAT changes measured with [$^{18}$F]FE-PE2I PET in patients with Parkinson's disease: a validation study,'' NeuroImage: Clinical, vol. 37, 103347, 2023. doi: 10.1016/j.nicl.2023.103347. PMID: 36822016

\bibitem{kim2022}
H.K.~Kim, et al., ``Temporal trajectory model for dopaminergic input to the striatal subregions in Parkinson's disease,'' \textit{Parkinsonism \& Related Disorders}, vol.~105, pp.~40--46, 2022. doi:~10.1016/j.parkreldis.2022.08.022.

\bibitem{kipf2017}
T.~N.~Kipf and M.~Welling, ``Semi-supervised classification with graph convolutional networks,'' in \textit{Proc.\ ICLR}, 2017.

\bibitem{kish1988}
Kish, S.J., et al. (1988).
\newblock Uneven pattern of dopamine loss in the striatum of patients with idiopathic Parkinson's disease.
\newblock \emph{New England Journal of Medicine}, 318(14), 876--880.

\bibitem{knowles2009}
T.~P.~J.~Knowles, C.~A.~Waudby, G.~L.~Devlin, S.~I.~A.~Cohen, A.~Aguzzi, M.~Vendruscolo, E.~M.~Terentjev, M.~E.~Welland, and C.~M.~Dobson, ``An analytical solution to the kinetics of breakable filament assembly,'' \textit{Science}, vol.~326, no.~5959, pp.~1533--1537, 2009.

\bibitem{koehler2008}
N.~K.~U.~Koehler \textit{et~al.}, ``A novel, highly sensitive and specific enzyme-linked immunosorbent assay (ELISA) for the measurement of total human alpha-synuclein in plasma and cerebrospinal fluid,'' \textit{Alzheimer's Dement.}, vol.~4, no.~4, Suppl., p.~T515, 2008.

% --- L ---

\bibitem{kong1994sequential} Kong, Augustine and Liu, Jun S. and Wong, Wing Hung ``Sequential imputations and Bayesian missing data problems,'' JASA, vol. 89, no. 425, pp. 278--288, 1994. doi: 10.1080/01621459.1994.10476469

\bibitem{koval2021admap}
I.~Koval \textit{et~al.}, ``AD Course Map charts Alzheimer's disease progression,'' \textit{Scientific Reports}, vol.~11, 8020, 2021. doi:~10.1038/s41598-021-87434-1.

\bibitem{kumar2020antibody}
S.~T.~Kumar \textit{et~al.}, ``How specific are the conformation-specific $\alpha$-synuclein antibodies? Characterization and validation of 16 $\alpha$-synuclein conformation-specific antibodies,'' \textit{Neurobiology of Disease}, vol.~146, p.~105086, 2020.

\bibitem{latourelle2017pd}
J.~C.~Latourelle \textit{et~al.}, ``Large-scale identification of clinical and genetic predictors of motor progression in patients with newly diagnosed Parkinson's disease,'' \textit{JAMA Neurol.}, vol.~74, no.~3, pp.~318--326, 2017.

\bibitem{laubenbacher2024}
R.~Laubenbacher, J.~P.~Sluka, and J.~A.~Glazier, ``Using digital twins in viral infection,'' \textit{Science}, vol.~383, no.~6681, pp.~303--305, 2024.

\bibitem{lee2018}
C.~Lee, W.~R.~Zame, J.~Yoon, and M.~van~der~Schaar, ``DeepHit: A deep learning approach to survival analysis with competing risks,'' in \textit{Proc.\ AAAI Conf.\ Artif.\ Intell.}, 2018.

\bibitem{lee2018deephit}
C.~Lee, W.~R.~Zame, J.~Yoon, and M.~van~der~Schaar, ``DeepHit: A deep learning approach to survival analysis with competing risks,'' in \textit{Proc.\ AAAI Conf.\ Artif.\ Intell.}, 2018.

\bibitem{lee2019}
C.~Lee, J.~Yoon, and M.~van~der~Schaar, ``Dynamic-DeepHit: A deep learning approach for dynamic survival analysis with competing risks based on longitudinal data,'' \textit{IEEE Trans.\ Biomed.\ Eng.}, vol.~67, no.~1, pp.~122--133, Jan. 2020.

\bibitem{lian2024}
J.~Lian, X.~Luo, C.~Shan, \textit{et~al.}, ``Personalized progression modelling and prediction in Parkinson's disease with a novel multi-modal graph approach,'' \textit{npj Parkinsons Dis.}, vol.~10, p.~229, 2024.

\bibitem{lian2024adamedgraph}
J.~Lian, X.~Luo, C.~Shan, \textit{et~al.}, ``Personalized progression modelling and prediction in Parkinson's disease with a novel multi-modal graph approach,'' \textit{npj Parkinsons Dis.}, vol.~10, p.~229, 2024.

\bibitem{little2019}
R.~J.~A.~Little and D.~B.~Rubin, \textit{Statistical Analysis with Missing Data}, 3rd~ed. Wiley, 2019.

% --- M ---

\bibitem{liu1998is}
J.~S.~Liu and R.~Chen, ``Sequential Monte Carlo methods for dynamic systems,'' \textit{Journal of the American Statistical Association}, vol.~93, no.~443, pp.~1032--1044, 1998.

% --- Added 2026-04-10: Post-SR mechanistic DT papers ---

\bibitem{liu2023gfap}
C.~Liu \textit{et~al.}, ``GFAP as an astroglial biomarker in neurodegenerative diseases: recent insights and pathophysiologic relevance,'' \textit{Journal of Neuroinflammation}, vol.~20, 132, 2023. doi:~10.1186/s12974-023-02812-y.

\bibitem{liu2025antibody}
Y.~Liu \textit{et~al.}, ``Comparative assessment of binding and functional differences between clinical-stage $\alpha$-synuclein antibodies,'' \textit{Biomedicine \& Pharmacotherapy}, vol.~183, p.~117870, 2025.

% --- Added 2026-04-09: IS-native convergence diagnostics + NUTS false-convergence methodology ---

\bibitem{maetzler2016}
W.~Maetzler, J.~Domingos, K.~Srulijes, J.~J.~Ferreira, and B.~R.~Bloem, ``Quantitative wearable sensors for objective assessment of Parkinson's disease,'' \textit{Mov.\ Disord.}, vol.~28, no.~12, pp.~1628--1637, 2013.

\bibitem{majbour2016}
N.~K.~Majbour, N.~N.~Vaikath, K.~D.~van~Dijk, M.~T.~Ardah, S.~Varghese, L.~B.~Vesterager, L.~P.~Montezinho, S.~Poole, B.~Safieh-Garabedian, T.~Tokuda, C.~E.~Teunissen, H.~W.~Berendse, W.~D.~J.~van~de~Berg, and O.~M.~A.~El-Agnaf, ``Oligomeric and phosphorylated alpha-synuclein as potential CSF biomarkers for Parkinson's disease,'' \textit{Molecular Neurodegeneration}, vol.~11, Art.\ no.~7, 2016, doi:10.1186/s13024-016-0072-9.

\bibitem{majbour2022}
N.~K.~Majbour, A.~J.~Aasly, L.~Abdi, et al., ``Disease-associated alpha-synuclein aggregates as biomarkers of Parkinson disease clinical stage,'' \textit{Neurology}, vol.~99, pp.~e2460--e2472, 2022. doi:~10.1212/WNL.0000000000201199.

% --- Added 2026-04-10: Missing citations flagged by literature validation audit ---

\bibitem{makarov2025}
V.~Makarov \textit{et~al.}, ``DT-GPT: Digital twin via GPT for clinical trajectory forecasting,'' \textit{npj Digit.\ Med.}, 2025.

\bibitem{mammana2024}
A.~Mammana, S.~Baiardi, A.~Rossi, S.~Orrù, and P.~Bhatt, ``LAG time is the most reliable kinetic parameter of the alpha-synuclein seed amplification assay,'' \textit{Clinical Chemistry and Laboratory Medicine}, vol.~62, pp.~1450--1460, 2024. doi:~10.1515/cclm-2023-1345.

\bibitem{mao2025nonmem} Mao, Bingyu and Gao, Yue and Xu, Christine and Macha, Sreeraj and Shao, Shuai and Ahamadi, Malidi ``Opportunities for AI-based Model-informed Drug Development: NONMEM vs AI-based Population PK,'' AAPS J, vol. 28, no. 1, pp. 21, 2025. doi: 10.1208/s12248-025-01121-x

\bibitem{mapie2022}
V.~Taquet, V.~Blot, T.~Morzadec, L.~Lacombe, and N.~Brunel, ``MAPIE: an open-source library for distribution-free uncertainty quantification,'' \textit{arXiv:2207.12274}, 2022.

\bibitem{marek2011}
K.~Marek \textit{et~al.}, ``The Parkinson Progression Marker Initiative (PPMI),'' \textit{Prog.\ Neurobiol.}, vol.~95, no.~4, pp.~629--635, 2011.

\bibitem{marek2018ppmi} Marek, Kenneth and Chowdhury, Sohini and Siderowf, Andrew and Lasch, Shirley and Coffey, Christopher S. and Caspell-Garcia, Chelsea and Simuni, Tanya and others ``The Parkinson's progression markers initiative (PPMI)---establishing a PD biomarker cohort,'' Ann Clin Transl Neurol, vol. 5, no. 12, pp. 1460--1477, 2018. doi: 10.1002/acn3.644

\bibitem{margossian2021nestedr}
C.~C.~Margossian, M.~D.~Hoffman, P.~Sountsov, L.~Riou-Durand, A.~Vehtari, and A.~Gelman, ``Nested $\widehat{R}$: Assessing the convergence of Markov chain Monte Carlo when running many short chains,'' \textit{Bayesian Analysis}, vol.~19, no.~1, pp.~289--321, 2024.

\bibitem{marras2018}
C.~Marras \textit{et~al.}, ``Prevalence of Parkinson's disease across North America,'' \textit{npj Parkinson's Dis.}, vol.~4, p.~21, 2018.

\bibitem{marshall2016goodpractices} Marshall, Shawn F. and Burghaus, Rolf and Cosson, Valerie and Cheung, S. Y. Amy and Chenel, Marylore and DellaPasqua, Oscar and others ``Good practices in model-informed drug discovery and development,'' CPT:PSP, vol. 5, no. 3, pp. 93--122, 2016. doi: 10.1002/psp4.12049

\bibitem{marshall2019midd} Marshall, Shawn and Madabushi, Rajanikanth and Manolis, Efthymios and others ``Model-informed drug discovery and development: Current industry good practice and regulatory expectations and future perspectives,'' CPT:PSP, vol. 8, no. 2, pp. 87--96, 2019. doi: 10.1002/psp4.12372

\bibitem{marshall2023harmonized} Marshall, Shawn and Ahamadi, M. and others ``Model-informed drug development: Steps toward harmonized guidance,'' Clin Pharmacol Ther, vol. 114, no. 5, pp. 962--965, 2023. doi: 10.1002/cpt.3006

\bibitem{martinez2015}
P.~Martinez-Martin, C.~Rodriguez-Blazquez, M.~J.~Forjaz, B.~Frades-Payo, W.~Ag{\"u}era-Ortiz, G.~T.~Weintraub, A.~Riesco, G.~Kurtis, and K.~R.~Chaudhuri, ``Neuropsychiatric symptoms and caregiver's burden in Parkinson's disease,'' \textit{Parkinsonism Relat.\ Disord.}, vol.~21, no.~6, pp.~629--634, 2015.

\bibitem{matsui2026stance}
K.~Matsui, Y.~Suzuki, C.~E.~Smith, T.~Nakamura, T.~Endo, S.~Sakoda, P.~G.~Morasso, and T.~Nomura, ``Digital twins of upright stance reveal mechanistic bifurcations underlying Parkinsonian sway phenotypes,'' \textit{bioRxiv preprint}, 2026. doi:~10.1101/2026.03.09.710685.

% --- Added 2026-04-10: Systematic review + HLME methodology ---

\bibitem{mattei2019}
P.-A.~Mattei and J.~Frisch, ``MIWAE: Deep generative modelling and imputation of incomplete data sets,'' in \textit{Proc.\ ICML}, 2019, pp.~4413--4423.

\bibitem{mattsson2019}
N.~Mattsson \textit{et~al.}, ``Association of plasma neurofilament light with neurodegeneration in patients with Alzheimer disease,'' \textit{JAMA Neurol.}, vol.~74, no.~5, pp.~557--566, 2017.

\bibitem{meisl2014}
G.~Meisl, J.~B.~Kirkegaard, P.~Arosio, T.~C.~T.~Michaels, M.~Vendruscolo, C.~M.~Dobson, S.~Linse, and T.~P.~J.~Knowles, ``Molecular mechanisms of protein aggregation from global fitting of kinetic models,'' \textit{Nature Protocols}, vol.~11, no.~2, pp.~252--272, 2016.

\bibitem{mekki2024}
O.~Mekki \textit{et~al.}, ``LSTM vs.\ Kolmogorov--Arnold Networks for MDS-UPDRS prediction,'' \textit{arXiv}, 2024; 2412.15089.

\bibitem{mindmend2025}
B.~Dupre, M.~Terry, J.~Shupe, and K.~House, ``Biomarker detection using layered receptor and electrode configuration,'' U.S.\ Patent Application US20250334570A1, Oct.~30, 2025.

% ============================================================================
% References added for Appendix D: Mechanistic Digital Twin Review
% ============================================================================

\bibitem{modrak2022}
M.~Modrak, et al., ``Simulation-based calibration checking for Bayesian computation: the choice of test quantities shapes sensitivity,'' \textit{Bayesian Analysis}, vol.~1, no.~1, pp.~1--28, 2022. doi:~10.1214/22-BA1404.

\bibitem{mollenhauer2008elisa}
B.~Mollenhauer \textit{et~al.}, ``Direct quantification of CSF $\alpha$-synuclein by ELISA and first cross-sectional study in patients with neurodegeneration,'' \textit{Experimental Neurology}, vol.~213, no.~2, pp.~315--325, 2008.

\bibitem{mollenhauer2011}
B.~Mollenhauer, J.~J.~Locascio, W.~Schulz-Schaeffer, F.~Sixel-D{\"o}ring, C.~Trenkwalder, and M.~G.~Schlossmacher, ``$\alpha$-Synuclein and tau concentrations in cerebrospinal fluid of patients presenting with parkinsonism: a cohort study,'' \textit{The Lancet Neurology}, vol.~10, no.~3, pp.~230--240, 2011, doi:10.1016/S1474-4422(11)70014-X.

\bibitem{mollenhauer2017}
B.~Mollenhauer, C.~J.~Caspell-Garcia, C.~T.~Coffey, \textit{et~al.}, ``Longitudinal CSF biomarkers in patients with early Parkinson disease and healthy controls,'' \textit{Neurology}, vol.~89, no.~19, pp.~1959--1969, 2017. doi:\,10.1212/WNL.0000000000004609.

\bibitem{mollenhauer2017csf}
B.~Mollenhauer, C.~J.~Caspell-Garcia, C.~T.~Coffey, \textit{et~al.}, ``Longitudinal CSF biomarkers in patients with early Parkinson disease and healthy controls,'' \textit{Neurology}, vol.~89, no.~19, pp.~1959--1969, 2017, doi:10.1212/WNL.0000000000004609.

\bibitem{mollenhauer2019longitudinal}
B.~Mollenhauer \textit{et~al.}, ``Longitudinal analyses of cerebrospinal fluid $\alpha$-synuclein in prodromal and early Parkinson's disease,'' \textit{Movement Disorders}, vol.~34, no.~9, pp.~1354--1364, 2019. doi:~10.1002/mds.27806.

\bibitem{moons2014}
K.~G.~M.~Moons, J.~A.~H.~de~Groot, W.~Bouwmeester, \textit{et~al.}, ``Critical appraisal and data extraction for systematic reviews of prediction modelling studies: the CHARMS checklist,'' \textit{PLoS Med.}, vol.~11, no.~10, Art.\ no.~e1001744, 2014.

% --- N ---

\bibitem{mould2013basic2} Mould, Diane R. and Upton, Richard N. ``Basic concepts in population modeling, simulation, and model-based drug development---Part 2: Introduction to pharmacokinetic modeling methods,'' CPT:PSP, vol. 2, pp. e38, 2013. doi: 10.1038/psp.2013.14

\bibitem{mould2014basic3} Mould, Diane R. and Upton, Richard N. ``Basic concepts in population modeling, simulation, and model-based drug development: Part 3---introduction to pharmacodynamic modeling methods,'' CPT:PSP, vol. 3, pp. e88, 2014. doi: 10.1038/psp.2013.71

\bibitem{musuamba2021credibility} Musuamba, Flora T. and others ``Scientific and regulatory evaluation of mechanistic in silico drug and disease models in drug development: Building model credibility,'' CPT:PSP, vol. 10, no. 8, pp. 804--825, 2021. doi: 10.1002/psp4.12669

\bibitem{naeini2015ece}
M.~P.~Naeini, G.~Cooper, and M.~Hauskrecht, ``Obtaining well calibrated probabilities using Bayesian binning,'' in \textit{Proc.\ AAAI Conf.\ Artif.\ Intell.}, 2015, pp.~2901--2907.

\bibitem{nair2026computationalpd} Nair, S.S. and Guha, A. and Chakravarthy, S. and Shaikh, A.G. ``Computational Modeling of PD Across Scales: Mechanisms to Biomarkers to Personalized Therapies,'' Brain Sciences, vol. 16, no. 2, pp. 175, 2026. doi: 10.3390/brainsci16020175

\bibitem{nakagawa2013}
S.~Nakagawa and H.~Schielzeth, ``A general and simple method for
obtaining $R^2$ from generalized linear mixed-effects models,''
\textit{Methods Ecol.\ Evol.}, vol.~4, no.~2, pp.~133--142, 2013.

\bibitem{NASEM2024DigitalTwins} National Academies of Sciences, Engineering, and Medicine ``Foundational Research Gaps and Future Directions for Digital Twins,'' 2024. doi: 10.17226/26894

\bibitem{nalls2019}
M.~A.~Nalls, C.~Blauwendraat, C.~L.~Vallerga, \textit{et~al.}, ``Identification of novel risk loci, causal insights, and heritable risk for Parkinson's disease: a meta-analysis of genome-wide association studies,'' \textit{Lancet Neurol.}, vol.~18, no.~12, pp.~1091--1102, 2019. doi:\,10.1016/S1474-4422(19)30320-5.

\bibitem{nestor2019temporal}
B.~Nestor \textit{et~al.}, ``Feature robustness in non-stationary health records: caveats to deployable model performance in common clinical machine learning tasks,'' in \textit{Proc.\ Mach.\ Learn.\ Healthcare (MLHC)}, 2019.

\bibitem{nilashi2022}
M.~Nilashi, O.~Ibrahim, H.~Ahmadi, \textit{et~al.}, ``Predicting Parkinson's disease progression: evaluation of ensemble methods,'' \textit{J.\ Healthc.\ Eng.}, vol.~2022, Art.\ no.~7634521, 2022.

% --- O ---

\bibitem{ogilvie2008}
D.~Ogilvie, D.~Fayter, M.~Petticrew, \textit{et~al.}, ``The harvest plot: a method for synthesising evidence about the differential effects of interventions,'' \textit{BMC Med.\ Res.\ Methodol.}, vol.~8, p.~8, 2008.

% --- P ---

\bibitem{ogle2020}
K.~Ogle, J.~J.~Barber, and M.~R.~Whitham, ``Ensuring identifiability in hierarchical mixed effects Bayesian models,'' \textit{Ecological Applications}, vol.~30, e02159, 2020. doi:~10.1002/eap.2159.

\bibitem{oh2012} M. Oh, J. S. Kim, J. Y. Kim, et al., ``Subregional patterns of preferential striatal dopamine transporter loss differ in Parkinson disease, progressive supranuclear palsy, and multiple-system atrophy,'' Journal of Nuclear Medicine, vol. 53, no. 3, pp. 399--406, 2012. doi: 10.2967/jnumed.111.095224

\bibitem{oliverassalva2013} A. Oliveras-Salv\'a, A. Van der Perren, N. Casadei, et al., ``rAAV2/7 vector-mediated overexpression of alpha-synuclein in mouse substantia nigra induces protein aggregation and progressive dose-dependent neurodegeneration,'' Molecular Neurodegeneration, vol. 8, p. 44, 2013. doi: 10.1186/1750-1326-8-44

\bibitem{orru2025}
S.~Orrù, C.~D.~Bhatt, T.~Bhatt, et al., ``Alpha-synuclein seed amplification assay time-to-threshold predicts cognitive decline in Parkinson's disease,'' \textit{The Lancet Neurology}, vol.~24, pp.~230--241, 2025. doi:~10.1016/S1474-4422(25)00015-3.

\bibitem{pagano2022}
G.~Pagano \textit{et~al.}, ``Trial of prasinezumab in early-stage Parkinson's disease,'' \textit{New Engl.~J.~Med.}, vol.~387, no.~5, pp.~421--432, 2022.

\bibitem{pagano2024}
G.~Pagano \textit{et~al.}, ``Sustained effect of prasinezumab on Parkinson's disease motor progression in the open-label extension of the PASADENA trial,'' \textit{Nature Medicine}, vol.~30, no.~12, pp.~3669--3675, 2024.

\bibitem{pagano2024subgroup}
G.~Pagano \textit{et~al.}, ``Prasinezumab slows motor progression in rapidly progressing early-stage Parkinson's disease,'' \textit{Nature Medicine}, vol.~30, no.~4, pp.~1096--1103, 2024.

\bibitem{page2021}
M.~J.~Page, J.~E.~McKenzie, P.~M.~Bossuyt, \textit{et~al.}, ``The PRISMA 2020 statement: an updated guideline for reporting systematic reviews,'' \textit{BMJ}, vol.~372, Art.\ no.~n71, 2021.

\bibitem{pandya2019}
S.~Pandya, et al., ``Predictive model of spread of Parkinson's pathology using network diffusion,'' \textit{NeuroImage}, vol.~192, pp.~178--194, 2019. doi:~10.1016/j.neuroimage.2019.02.058.

\bibitem{paper1}
B.~Dupre, ``NSD-ISS stage prediction with calibrated uncertainty for Parkinson's disease (this dissertation, Chapter~3),'' \textit{IEEE J.\ Biomed.\ Health Inform.}, 2026.

\bibitem{paper2}
B.~Dupre, ``Stage-conditioned graph-informed imputation with heteroscedastic uncertainty for Parkinson's disease clinical data (this dissertation, Chapter~4),'' \textit{IEEE J.\ Biomed.\ Health Inform.}, 2026.

\bibitem{paper3}
B.~Dupre, ``Graph-informed digital twins for predicting NSD-ISS stage transitions in Parkinson's disease (this dissertation, Chapter~5),'' \textit{IEEE J.\ Biomed.\ Health Inform.}, 2026.

\bibitem{parkkinen2011}
L.~Parkkinen \textit{et~al.}, ``Does levodopa accelerate the pathologic
process in Parkinson disease brain?'' \textit{Neurology}, vol.~77,
no.~15, pp.~1420--1426, 2011.

\bibitem{pasquini2019}
J.~Pasquini, et al., ``Clinical implications of early caudate dysfunction in Parkinson's disease,'' \textit{Journal of Neurology, Neurosurgery, and Psychiatry}, vol.~90, no.~6, pp.~681--687, 2019. doi:~10.1136/jnnp-2018-320157.

\bibitem{patterson2019}
J.R.~Patterson, et al., ``Time course and magnitude of alpha-synuclein inclusion formation and nigrostriatal degeneration in the rat model of synucleinopathy triggered by intrastriatal alpha-synuclein preformed fibrils,'' \textit{Neurobiology of Disease}, vol.~121, pp.~1--12, 2019. doi:~10.1016/j.nbd.2018.09.001.

\bibitem{petricca2022}
L.~Petricca, N.~Chiki, M.~F.~Bhatt, \textit{et~al.}, ``Comparative analysis of total alpha-synuclein ($\alpha$SYN) immunoassays reveals that they do not capture the diversity of modified $\alpha$SYN proteoforms,'' \textit{J.~Parkinsons Dis.}, vol.~12, no.~5, pp.~1449--1462, 2022, doi:10.3233/JPD-223285.

\bibitem{postuma2015}
R.~B.~Postuma, D.~Berg, M.~Stern, \textit{et~al.}, ``MDS clinical diagnostic criteria for Parkinson's disease,'' \textit{Mov.\ Disord.}, vol.~30, no.~12, pp.~1591--1601, 2015.

% --- R ---

\bibitem{powell2018}
F.~Powell, M.~Tosun, R.~Sadeghi, P.~Weiner, and A.~Raj, ``Preserved structural network organization mediates pathology spread in Alzheimer's disease despite loss of white matter tract integrity,'' \textit{Journal of Alzheimer's Disease}, vol.~65, pp.~747--764, 2018. doi:~10.3233/JAD-180515.

% === Paper 8a identifiability additions (2026-04-11) ===

\bibitem{ppmi2011}
Parkinson Progression Marker Initiative (2011).
\newblock The Parkinson Progression Marker Initiative (PPMI).
\newblock \emph{Progress in Neurobiology}, 95(4), 629--635.

\bibitem{prokhorenkova2018catboost}
L.~Prokhorenkova, G.~Gusev, A.~Vorobev, A.~V.~Dorogush, and A.~Gulin, ``CatBoost: unbiased boosting with categorical features,'' in \textit{Proc.\ NeurIPS}, 2018.

\bibitem{rackauckas2017}
C.~Rackauckas and Q.~Nie, ``DifferentialEquations.jl---A performant and feature-rich ecosystem for solving differential equations in Julia,'' \textit{J.\ Open Research Software}, vol.~5, no.~1, p.~15, 2017.

\bibitem{rackauckas2020}
C.~Rackauckas, Y.~Ma, J.~Martensen, \textit{et~al.}, ``Universal differential equations for scientific machine learning,'' \textit{arXiv preprint arXiv:2001.04385}, 2020.

\bibitem{rackauckas2020universal} Rackauckas, Christopher and Ma, Yingbo and Martensen, Julius and Warner, Collin and Zubov, Kirill and Supekar, Rohit and Skinner, Dominic and Ramadhan, Ali and Edelman, Alan ``Universal Differential Equations for Scientific Machine Learning,'' 2020. arXiv:2001.04385

\bibitem{raj2012}
A.~Raj, A.~Kuceyeski, and M.~Weiner, ``A network diffusion model of disease progression in dementia,'' \textit{Neuron}, vol.~73, no.~6, pp.~1204--1215, 2012.

\bibitem{raj2019}
A.~Raj, E.~Powell, \textit{et~al.}, ``Predictive model of spread of Parkinson's pathology using network diffusion,'' \textit{NeuroImage}, vol.~192, pp.~178--194, 2019.

\bibitem{rajkomar2018}
A.~Rajkomar, J.~Dean, and I.~Kohane, ``Machine learning in medicine,'' \textit{N.\ Engl.\ J.\ Med.}, vol.~380, no.~14, pp.~1347--1358, 2019.

\bibitem{rajkomar2018ensuring}
A.~Rajkomar, M.~Hardt, M.~D.~Howell, G.~Corrado, and M.~H.~Chin, ``Ensuring fairness in machine learning to advance health equity,'' \textit{Ann.\ Intern.\ Med.}, vol.~169, no.~12, pp.~866--872, 2018.

\bibitem{rastegar2019}
D.~Ahmadi~Rastegar, N.~Ho, G.~M.~Halliday, and N.~Dzamko, ``Parkinson's progression prediction using machine learning and serum cytokines,'' \textit{npj Parkinsons Dis.}, vol.~5, p.~2, 2019.

\bibitem{raue2009}
A.~Raue, C.~Kreutz, T.~Maiwald, J.~Bachmann, M.~Schilling, U.~Klingm{\"u}ller, and J.~Timmer, ``Structural and practical identifiability analysis of partially observed dynamical models by exploiting the profile likelihood,'' \textit{Bioinformatics}, vol.~25, no.~15, pp.~1923--1929, 2009, doi:10.1093/bioinformatics/btp358.

\bibitem{reconsider2025nsd}
\textit{et~al.}, ``Reconsidering the clinical foundations of NSD-ISS in PD staging,'' \textit{Mov.\ Disord.}, 2025. doi:\,10.1002/mds.70061.

\bibitem{ren2021}
X.~Ren, J.~Lin, G.~T.~Stebbins, C.~G.~Goetz, and S.~Luo, ``Prognostic modeling of Parkinson's disease progression using early longitudinal patterns of change,'' \textit{Mov.\ Disord.}, vol.~36, no.~12, pp.~2808--2816, 2021.

\bibitem{ren2024}
J.~Ren, H.~Xie, Y.~Weng, \textit{et~al.}, ``Longitudinal decline in DAT binding in Parkinson's disease: connections with sleep disturbances,'' \textit{BMC Medicine}, vol.~22, Art.\ no.~520, 2024.

\bibitem{ren2024datsleep} Ren, Junli and Xie, Haobo and Weng, Yiyun and others ``Longitudinal decline in DAT binding in PD: connections with sleep disturbances,'' BMC Medicine, vol. 22, pp. 550, 2024. doi: 10.1186/s12916-024-03766-5

\bibitem{ren2025}
J.~Ren, L.~Bhatt, G.~Guo, and A.~Raj, ``Connectome-based biophysical models of pathological protein spreading in neurodegenerative diseases,'' \textit{PLoS Comput.\ Biol.}, vol.~21, no.~1, e1012743, 2025, doi:10.1371/journal.pcbi.1012743.

\bibitem{ribba2024}
B.~Ribba, et al., ``Disease progression modeling of MDS-UPDRS in Parkinson's disease using PPMI data,'' \textit{Journal of Parkinson's Disease}, 2024.

\bibitem{righetti2025}
E.~Righetti, L.~Marchetti, E.~Domenici, and F.~Reali, ``A mechanistic model of pure and lipidic $\alpha$-synuclein aggregation for advancing Parkinson's therapies,'' \textit{Communications Chemistry}, vol.~8, art.~186, 2025.

\bibitem{riley2020}
R.~D.~Riley \textit{et~al.}, ``Calculating the sample size required for developing a clinical prediction model,'' \textit{BMJ}, vol.~368, Art.\ no.~m441, 2020.

\bibitem{rissanen2022}
E.~Rissanen \textit{et~al.}, ``Machine learning applications for Parkinson's disease: a systematic review,'' \textit{NPJ Digit.\ Med.}, 2022.

\bibitem{rizou2025}
V.~Rizou, N.~Grammalidis, P.~Daras, and K.~Dimitropoulos, ``PDualNet: a deep learning framework for joint prediction of Parkinson's disease progression subtype and MDS-UPDRS scores,'' \textit{Sci.\ Rep.}, vol.~15, p.~41931, 2025.

\bibitem{roberts2017temporal}
D.~R.~Roberts \textit{et~al.}, ``Cross-validation strategies for data with temporal, spatial, hierarchical, or phylogenetic structure,'' \textit{Ecography}, vol.~40, no.~8, pp.~913--929, 2017.

\bibitem{roy2019convergence}
V.~Roy, ``Convergence diagnostics for Markov chain Monte Carlo,'' \textit{Annual Review of Statistics and Its Application}, vol.~6, pp.~387--412, 2019.

\bibitem{rubin1976}
D.~B.~Rubin, ``Inference and missing data,'' \textit{Biometrika}, vol.~63, no.~3, pp.~581--592, 1976.

\bibitem{rukavina2025}
K.~Rukavina \textit{et~al.}, ``DaT-SPECT uptake predicts subsequent
levodopa equivalent daily dose requirements in Parkinson's disease,''
2025.

\bibitem{russo2025}
M.~J.~Russo \textit{et~al.}, ``Application of the NSD-ISS biological staging framework in the BioFIND cohort,'' \textit{npj Parkinson's Dis.}, 2025.

% --- S ---

\bibitem{rutledge2024}
J.~Rutledge, H.~Lehallier, G.~K.~Zarifkar, et al., ``Comprehensive proteomics of CSF, plasma, and urine identify DDC and other biomarkers of early Parkinson's disease,'' \textit{Acta Neuropathologica}, vol.~147, art.~71, 2024. doi:~10.1007/s00401-024-02720-0.

\bibitem{sadaei2022}
H.~J.~Sadaei, A.~Cordova-Palomera, J.~Lee, \textit{et~al.}, ``Genetically-informed prediction of short-term Parkinson's disease progression,'' \textit{npj Parkinsons Dis.}, vol.~8, p.~143, 2022.

\bibitem{salmanpour2022}
M.~R.~Salmanpour, M.~Shamsaei, G.~Hajianfar, \textit{et~al.}, ``Longitudinal clustering analysis and prediction of Parkinson's disease progression using radiomics and hybrid machine learning,'' \textit{Quant.\ Imaging Med.\ Surg.}, vol.~12, no.~2, pp.~906--919, 2022.

\bibitem{sampedro2020nfl}
F.~Sampedro \textit{et~al.}, ``Serum neurofilament light chain levels reflect cortical neurodegeneration in de novo Parkinson's disease,'' \textit{Parkinsonism \& Related Disorders}, vol.~74, pp.~43--49, 2020. doi:~10.1016/j.parkreldis.2020.04.009.

\bibitem{sarkar2026}
S.~Sarkar and M.~Podder, ``Profiling the Braak progression in Parkinson's disease: transcriptomics and ML-driven identification of progressive biomarkers,'' \textit{Neuroscience}, vol.~547, p.~89, 2026.

\bibitem{savickarlsson2009aaps}
R.~M.~Savic and M.~O.~Karlsson, ``Importance of shrinkage in empirical Bayes estimates for diagnostics: problems and solutions,'' \textit{The AAPS Journal}, vol.~11, no.~3, pp.~558--569, 2009. doi:~10.1208/s12248-009-9133-0.

\bibitem{schafer2021}
A.~Schafer, et al., ``Bayesian physics-based modeling of tau propagation in Alzheimer's disease,'' \textit{Frontiers in Physiology}, vol.~12, 702975, 2021. doi:~10.3389/fphys.2021.702975.

\bibitem{schenk2017}
D.~B.~Schenk \textit{et~al.}, ``First-in-human assessment of PRX002, an anti-$\alpha$-synuclein monoclonal antibody, in healthy volunteers,'' \textit{Movement Disorders}, vol.~32, no.~2, pp.~211--218, 2017.

\bibitem{schmitzSteinkruger2021age}
H.~Schmitz-Steinkr\"{u}ger, C.~Lange, I.~Apostolova, F.~L.~Mathies, L.~Frings, S.~Klutmann, S.~Hellwig, P.~T.~Meyer, and R.~Buchert, ``Impact of age and sex correction on the diagnostic performance of dopamine transporter SPECT,'' \textit{Eur.\ J.\ Nucl.\ Med.\ Mol.\ Imaging}, vol.~48, no.~5, pp.~1445--1459, May 2021. doi:\,10.1007/s00259-020-05085-2.

\bibitem{schunck2025}
C.~Schunck, J.~Hemmerich, and G.~Nickel, ``Hierarchical approaches for integrating sparse, multivariate toxicological effect data in whole organism molecular dynamics,'' \textit{bioRxiv preprint}, 2025. doi:~10.1101/2025.02.15.638451.

\bibitem{seibyl1997jnm}
J.~P.~Seibyl \textit{et~al.}, ``Reproducibility of iodine-123-$\beta$-CIT SPECT brain measurement of dopamine transporters,'' \textit{Journal of Nuclear Medicine}, vol.~38, no.~9, pp.~1453--1459, 1997. PMID: 9293807.

\bibitem{seibyl2018dat}
J.~P.~Seibyl, ``DaT-SPECT in Parkinson's disease: clinical utility and limitations,'' \textit{J.\ Nucl.\ Med.}, 2018.

\bibitem{severson2021}
K.A.~Severson, et al., ``Discovery of Parkinson's disease states and disease progression modelling: a longitudinal data study using machine learning,'' \textit{The Lancet Digital Health}, vol.~3, no.~9, pp.~e555--e564, 2021. doi:~10.1016/S2589-7500(21)00101-1.

\bibitem{shafer2008}
G.~Shafer and V.~Vovk, ``A tutorial on conformal prediction,'' \textit{J.\ Mach.\ Learn.\ Res.}, vol.~9, pp.~371--421, 2008.

\bibitem{shahnawaz2023}
M.~Shahnawaz, A.~Mukherjee, S.~Pritzkow, \textit{et~al.}, ``Propagative $\alpha$-synuclein seeds as serum biomarkers for synucleinopathies,'' \textit{Nature Medicine}, vol.~29, no.~6, pp.~1448--1455, 2023.

\bibitem{shang2019}
J.~Shang, T.~Ma, C.~Xiao, and J.~Sun, ``Pre-training of graph augmented transformers for medication recommendation,'' in \textit{Proc.\ IJCAI}, 2019.

\bibitem{sheiner1980}
L.~B.~Sheiner and S.~L.~Beal, ``Evaluation of methods for estimating population pharmacokinetic parameters. I. Michaelis-Menten model: routine clinical pharmacokinetic data,'' \textit{J.\ Pharmacokinet.\ Biopharm.}, vol.~8, no.~6, pp.~553--571, 1980.

\bibitem{shwartzziv2022}
R.~Shwartz-Ziv and A.~Armon, ``Tabular data: Deep learning is not all you need,'' \textit{Information Fusion}, vol.~81, pp.~84--90, 2022.

\bibitem{siderowf2020}
A.~Siderowf, D.~Jennings, M.~Stern, \textit{et~al.}, ``Clinical and imaging progression in the PARS cohort: long-term follow-up,'' \textit{Mov.\ Disord.}, vol.~35, no.~9, pp.~1550--1557, 2020.

\bibitem{simon2016}
N.~Simon, et al., ``A combined pharmacokinetic/pharmacodynamic model of levodopa motor response and dyskinesia in Parkinson's disease patients,'' \textit{European Journal of Clinical Pharmacology}, vol.~72, pp.~423--434, 2016.

\bibitem{simpson2024}
M.J.~Simpson, et al., ``Making predictions using poorly identified mathematical models,'' \textit{Bulletin of Mathematical Biology}, vol.~86, no.~7, 83, 2024. doi:~10.1007/s11538-024-01294-0.

\bibitem{simuni2018} T. Simuni, C. L. Siderowf, S. Lasch, et al., ``Longitudinal change of clinical and biological measures in early Parkinson's disease: PPMI cohort,'' Movement Disorders, vol. 33, no. 5, pp. 771--782, 2018. doi: 10.1002/mds.27361

\bibitem{simuni2020lrrk2}
T.~Simuni, B.~Uribe, J.~Y.~Cho, \textit{et~al.}, ``Clinical and dopamine transporter imaging characteristics of leucine rich repeat kinase 2 (LRRK2) and glucosylceramidase beta (GBA) Parkinson's disease participants in the Parkinson's Progression Markers Initiative,'' \textit{Mov.\ Disord.}, vol.~35, no.~5, pp.~833--844, 2020.

\bibitem{simuni2024}
T.~Simuni \textit{et~al.}, ``A biological definition of neuronal alpha-synuclein disease: towards an integrated staging system for research,'' \textit{Lancet Neurol.}, vol.~23, no.~2, pp.~178--190, Feb. 2024.

\bibitem{simuni2024nsd}
T.~Simuni \textit{et~al.}, ``A biological definition of neuronal alpha-synuclein disease: towards an integrated staging system for research,'' \textit{Lancet Neurol.}, vol.~23, no.~2, pp.~178--190, Feb. 2024.

\bibitem{simuni2024nsdiss}
T.~Simuni \textit{et~al.}, ``A biological definition of neuronal alpha-synuclein disease: towards an integrated staging system for research,'' \textit{Lancet Neurol.}, vol.~23, no.~2, pp.~178--190, Feb. 2024.

\bibitem{simuni2024val}
T.~Simuni \textit{et~al.}, ``Validation of the NSD-ISS in PPMI, PASADENA, and SPARK,'' \textit{npj Parkinson's Dis.}, 2024.

\bibitem{simuni2025}
T.~Simuni \textit{et~al.}, ``Longitudinal clinical and biomarker characteristics of non-manifesting carriers and participants with Lewy body spectrum disease in the PPMI study using the NSD-ISS: A prospective cohort study,'' \textit{Lancet Neurol.}, vol.~24, no.~2, pp.~127--142, Feb. 2025.

\bibitem{simuni2025progression}
T.~Simuni \textit{et~al.}, ``Longitudinal clinical and biomarker characteristics of non-manifesting carriers and participants with Lewy body spectrum disease in the PPMI study using the NSD-ISS: A prospective cohort study,'' \textit{Lancet Neurol.}, vol.~24, no.~2, pp.~127--142, Feb. 2025.

\bibitem{simuni2025reply}
T.~Simuni \textit{et~al.}, ``Reply to: Refutation of the $\alpha$Syn-SAA-based staging for Parkinson's progression,'' \textit{Mov.\ Disord.}, 2025. doi:\,10.1002/mds.30272.

\bibitem{south2019smc} South, Leah F. and Pettitt, Anthony N. and Drovandi, Christopher C. ``Sequential Monte Carlo samplers with independent Markov chain Monte Carlo proposals,'' Bayesian Analysis, vol. 14, no. 3, pp. 753--776, 2019. doi: 10.1214/18-ba1129

\bibitem{sreenivasan2025}
A.~P.~Sreenivasan, A.~Vaivade, Y.~Noui, \textit{et~al.}, ``Conformal prediction enables disease course prediction and allows individualized diagnostic uncertainty in multiple sclerosis,'' \textit{npj Digit.\ Med.}, vol.~8, Art.\ no.~224, 2025, doi:10.1038/s41746-025-01616-z.

\bibitem{stekhoven2012}
D.~J.~Stekhoven and P.~B{\"u}hlmann, ``MissForest---non-parametric missing value imputation for mixed-type data,'' \textit{Bioinformatics}, vol.~28, no.~1, pp.~112--118, Jan. 2012.

% --- T ---

\bibitem{sung2016}
C.~Sung, et al., ``Longitudinal decline of striatal subregional [18F]FP-CIT uptake in Parkinson's disease,'' \textit{Nuclear Medicine and Molecular Imaging}, vol.~50, no.~4, pp.~314--320, 2016. doi:~10.1007/s13139-016-0436-4.

\bibitem{talts2018}
S.~Talts, et al., ``Validating Bayesian inference algorithms with simulation-based calibration,'' \textit{arXiv preprint arXiv:1804.06788}, 2018.

% ============================================================================
% Phase 4 PK/PD defensive citations (added 2026-04-12)
% ============================================================================

\bibitem{temple2010}
R.~Temple, ``Enrichment of clinical study populations,'' \textit{Clin.\ Pharmacol.\ Ther.}, vol.~88, no.~6, pp.~774--778, 2010.

\bibitem{tokuda2010}
T.~Tokuda \textit{et~al.}, ``Detection of elevated levels of $\alpha$-synuclein oligomers in CSF from patients with Parkinson disease,'' \textit{Neurology}, vol.~75, no.~20, pp.~1766--1772, 2010.

\bibitem{tolosa2021}
E.~Tolosa, A.~Garrido, S.~W.~Scholz, and W.~Poewe, ``Challenges in the diagnosis of Parkinson's disease,'' \textit{Lancet Neurol.}, vol.~20, no.~5, pp.~385--397, 2021.

\bibitem{tomlinson2010ledd} Tomlinson, Claire L. and Stowe, Rebecca and Patel, Smitaa and Rick, Caroline and Gray, Richard and Clarke, Carl E. ``Systematic review of levodopa dose equivalency reporting in Parkinson's disease,'' Mov Disord, vol. 25, no. 15, pp. 2649--2653, 2010. doi: 10.1002/mds.23429

\bibitem{topol2019}
E.~J.~Topol, ``High-performance medicine: the convergence of human and artificial intelligence,'' \textit{Nat.\ Med.}, vol.~25, no.~1, pp.~44--56, 2019.

\bibitem{tossicibolt2017}
L.~Tossici-Bolt \textit{et~al.}, ``[$^{123}$I]FP-CIT ENC-DAT normal database: the impact of the reconstruction and quantification methods,'' \textit{EJNMMI Phys.}, vol.~4, Art.\ no.~8, 2017.

\bibitem{transtrum2015}
M.~K.~Transtrum, B.~B.~Machta, K.~S.~Brown, B.~C.~Daniels, C.~R.~Myers, and J.~P.~Sethna, ``Perspective: Sloppiness and emergent theories in physics, biology, and beyond,'' \textit{The Journal of Chemical Physics}, vol.~143, no.~1, p.~010901, 2015, doi:10.1063/1.4923066.

% --- U ---

\bibitem{ursino2020}
M.~Ursino, et al., ``Mathematical modeling and parameter estimation of levodopa motor response in patients with Parkinson disease,'' \textit{PLoS ONE}, vol.~15, no.~3, e0229729, 2020. doi:~10.1371/journal.pone.0229729.

\bibitem{vankessel2019}
S.~P.~van~Kessel, A.~K.~Frye, A.~O.~El-Gendy, \textit{et~al.}, ``Gut bacterial tyrosine decarboxylases restrict levels of levodopa in the treatment of Parkinson's disease,'' \textit{Nature Communications}, vol.~10, no.~1, Art.\ no.~310, 2019.

\bibitem{varrone2013ageadjusted}
A.~Varrone \textit{et~al.}, ``European multicentre database of healthy controls for [123I]FP-CIT SPECT (ENC-DAT): age-related effects, gender differences and evaluation of different methods of analysis,'' \textit{Eur.\ J.\ Nucl.\ Med.\ Mol.\ Imaging}, vol.~40, no.~2, pp.~213--227, 2013. doi:~10.1007/s00259-012-2276-8.

\bibitem{vehtari2012survey} Vehtari, Aki and Ojanen, Janne ``A survey of Bayesian predictive methods for model assessment, selection and comparison,'' Statistics Surveys, vol. 6, pp. 142--228, 2012. doi: 10.1214/12-SS102

\bibitem{vehtari2017practical} Vehtari, Aki and Gelman, Andrew and Gabry, Jonah ``Practical Bayesian model evaluation using leave-one-out cross-validation and WAIC,'' Stat Comput, vol. 27, no. 5, pp. 1413--1432, 2017. doi: 10.1007/s11222-016-9696-4

\bibitem{vehtari2017psis}
A.~Vehtari, A.~Gelman, and J.~Gabry, ``Practical Bayesian model evaluation using leave-one-out cross-validation and WAIC,'' \textit{Statistics and Computing}, vol.~27, no.~5, pp.~1413--1432, 2017, doi:10.1007/s11222-016-9696-4.

\bibitem{vehtari2021rhat}
A.~Vehtari, A.~Gelman, D.~Simpson, B.~Carpenter, and P.-C.~B{\"u}rkner, ``Rank-normalization, folding, and localization: An improved $\widehat{R}$ for assessing convergence of MCMC (with discussion),'' \textit{Bayesian Analysis}, vol.~16, no.~2, pp.~667--718, 2021.

\bibitem{vehtari2024psis} Vehtari, Aki and Simpson, Daniel and Gelman, Andrew and Yao, Yuling and Gabry, Jonah ``Pareto smoothed importance sampling,'' JMLR 25(72):1--58, 2024. arXiv:1507.02646

\bibitem{vela2022temporal}
D.~Vela \textit{et~al.}, ``Temporal quality degradation in AI models,'' \textit{Sci.\ Rep.}, vol.~12, no.~1, p.~11654, 2022.

\bibitem{velickovic2018}
P.~Veli\v{c}kovi\'{c}, G.~Cucurull, A.~Casanova, A.~Romero, P.~Li\`{o}, and Y.~Bengio, ``Graph attention networks,'' in \textit{Proc.\ ICLR}, 2018.

\bibitem{velickovic2018gat}
P.~Veli\v{c}kovi\'{c}, G.~Cucurull, A.~Casanova, A.~Romero, P.~Li\`{o}, and Y.~Bengio, ``Graph attention networks,'' in \textit{Proc.\ ICLR}, 2018.

\bibitem{venuto2023}
C.~S.~Venuto, G.~Smith, K.~Herbst, \textit{et~al.}, ``Predicting ambulatory capacity in Parkinson's disease to analyze progression, biomarkers, and trial design,'' \textit{Mov.\ Disord.}, vol.~38, no.~10, pp.~1774--1785, 2023.

\bibitem{veronneau2020}
F.~V\'{e}ronneau-Veilleux, et al., ``An integrative model of Parkinson's disease treatment including levodopa pharmacokinetics, dopamine kinetics, basal ganglia neurotransmission and motor action throughout disease progression,'' \textit{Journal of Pharmacokinetics and Pharmacodynamics}, vol.~47, pp.~539--562, 2020. doi:~10.1007/s10928-020-09719-w.

\bibitem{veronneau2020chaos} V\'eronneau-Veilleux, Florence and Ursino, Mauro and Robaey, Philippe and L\'evesque, Daniel and Nekka, Fahima ``Nonlinear pharmacodynamics of levodopa through Parkinson's disease progression,'' Chaos, vol. 30, no. 9, pp. 093146, 2020. doi: 10.1063/5.0014800

\bibitem{veronneau2020integrative} V\'eronneau-Veilleux, Florence and Robaey, Philippe and Ursino, Mauro and Nekka, Fahima ``An integrative model of PD treatment including levodopa pharmacokinetics, dopamine kinetics, basal ganglia neurotransmission and motor action,'' J Pharmacokinet Pharmacodyn, vol. 48, pp. 133--148, 2021. doi: 10.1007/s10928-020-09723-y

\bibitem{verschuur2019}
C.V.M.~Verschuur, et al., ``Randomized delayed-start trial of levodopa in Parkinson's disease,'' \textit{New England Journal of Medicine}, vol.~380, no.~4, pp.~315--324, 2019. doi:~10.1056/NEJMoa1809983.

% =============================================================================
% Phase 5 Task 5 + Task 6 + Task 7 citations (Papers 8a, 8b, 9, 10 chapters)
% Merged 2026-04-13 from phase5_literature_bibliography.bib
% =============================================================================

\bibitem{viceconti2021insilico} Viceconti, Marco and Pappalardo, Francesco and Rodriguez, Blanca and Horner, Marc and Bischoff, Jeff and Musuamba Tshinanu, Flora ``In silico trials: Verification, validation and uncertainty quantification of predictive models used in the regulatory evaluation of biomedical products,'' Methods, vol. 185, pp. 120--127, 2021. doi: 10.1016/j.ymeth.2020.01.011

\bibitem{viceconti2025credibilityml} Viceconti, Marco and Lanubile, Filippo and others ``Position Paper: Extending Credibility Assessment of In Silico Medicine Predictors to Machine Learning Predictors,'' IEEE J Biomed Health Inform, vol. 29, no. 7, pp. 5284--5290, 2025. doi: 10.1109/JBHI.2025.3552320

\bibitem{villaverde2016}
A.~F.~Villaverde, A.~Barreiro, and A.~Papachristodoulou, ``Structural identifiability of dynamic systems biology models,'' \textit{PLoS Computational Biology}, vol.~12, no.~10, p.~e1005153, 2016, doi:10.1371/journal.pcbi.1005153.

\bibitem{villaverde2022}
A.~F.~Villaverde, D.~Pathirana, F.~Fr{\"o}hlich, and J.~Hasenauer, ``Bayesian parameter estimation for dynamical models in systems biology,'' \textit{PLOS Computational Biology}, vol.~18, no.~10, Art.\ no.~e1010651, 2022.

\bibitem{vogel2023}
Vogel, J.W., et al. (2023).
\newblock Connectome-based modelling of neurodegenerative diseases: towards precision medicine and mechanistic insight.
\newblock \emph{Nature Reviews Neuroscience}, 24(10), 620--639.

\bibitem{vogel2024} J. W. Vogel, J. L. Young, S. Sch\"oll, R. La Joie, et al., ``Four distinct trajectories of tau deposition identified in Alzheimer's disease,'' Nature Medicine, vol. 27, no. 5, pp. 871--881, 2021. doi: 10.1038/s41591-021-01309-6

\bibitem{vogt2022}
A.~Vogt, ``Digital twins for Parkinson's disease medication dosing: Concept and challenges,'' 2022.

\bibitem{vovk2022}
V.~Vovk, A.~Gammerman, and G.~Shafer, \textit{Algorithmic Learning in a Random World}, 2nd~ed. Springer, 2022.

\bibitem{vovk2022algorithmic}
V.~Vovk, A.~Gammerman, and G.~Shafer, \textit{Algorithmic Learning in a Random World}, 2nd~ed. Springer, 2022.

\bibitem{vovk2022conformal}
V.~Vovk, A.~Gammerman, and G.~Shafer, \textit{Algorithmic Learning in a Random World}, 2nd~ed. Springer, 2022.

% --- W ---

\bibitem{vuong1989}
Q.H.~Vuong, ``Likelihood ratio tests for model selection and non-nested
hypotheses,'' \textit{Econometrica}, vol.~57, no.~2, pp.~307--333,
1989.

\bibitem{wakasugi2024combat} Wakasugi, Noritaka and Takano, Harumasa and Abe, Mitsunari and others ``Harmonizing multisite data with the ComBat method for enhanced PD diagnosis via DAT-SPECT,'' Front Neurol, vol. 15, pp. 1306546, 2024. doi: 10.3389/fneur.2024.1306546

\bibitem{wang2025}
S.~Wang \textit{et~al.}, ``Conditional neural ordinary differential equations for brain morphometry dynamics in Parkinson's disease,'' \textit{arXiv}, 2025; 2410.23429.

\bibitem{watts2024}
A.~Watts \textit{et~al.}, ``Deep learning for medication dosage trajectory prediction in Parkinson's disease,'' \textit{Sci.\ Rep.}, vol.~14, p.~7891, 2024.

\bibitem{weickenmeier2019}
J.~Weickenmeier \textit{et~al.}, ``A physics-based model explains the prion-like features of neurodegeneration in Alzheimer's disease, Parkinson's disease, and amyotrophic lateral sclerosis,'' \textit{Journal of the Mechanics and Physics of Solids}, vol.~124, pp.~264--281, 2019. doi:~10.1016/j.jmps.2018.10.013.

\bibitem{weihofen2019}
A.~Weihofen \textit{et~al.}, ``Development of an aggregate-selective, human-derived $\alpha$-synuclein antibody BIIB054 that ameliorates disease phenotypes in Parkinson's disease models,'' \textit{Neurobiology of Disease}, vol.~124, pp.~276--288, 2019.

\bibitem{wieland2021}
F.-G.~Wieland, A.~L.~Hauber, M.~Rosenblatt, C.~T\"onsing, and J.~Timmer, ``On structural and practical identifiability,'' \textit{Current Opinion in Systems Biology}, vol.~25, pp.~60--69, 2021.

\bibitem{winner2011}
B.~Winner, R.~Jappelli, S.~K.~Maji, \textit{et~al.}, ``In vivo demonstration that $\alpha$-synuclein oligomers are toxic,'' \textit{Proc.\ Natl.\ Acad.\ Sci.\ USA}, vol.~108, no.~10, pp.~4194--4199, 2011.

\bibitem{wolff2019}
R.~F.~Wolff, K.~G.~M.~Moons, R.~D.~Riley, \textit{et~al.}, ``PROBAST: a tool to assess the risk of bias and applicability of prediction model studies,'' \textit{Ann.\ Intern.\ Med.}, vol.~170, no.~1, pp.~51--58, 2019.

\bibitem{wouters2020}
O.~J.~Wouters \textit{et~al.}, ``Estimated research and development investment needed to bring a new medicine to market,'' \textit{JAMA}, vol.~323, no.~9, pp.~844--853, 2020.

\bibitem{wu2007psi}
D.~Wu, ``Measuring performance degradation of classifier due to non-stationarity in production environments,'' in \textit{Proc.\ ICDM Workshops}, pp.~727--732, 2007.

% --- X ---

\bibitem{xing2021}
Y.~Xing, A.~Sapuan, R.~A.~Dineen, and D.~P.~Auer, ``The spatiotemporal changes in dopamine, neuromelanin and iron characterizing Parkinson's disease,'' \textit{Brain}, vol.~144, no.~10, pp.~3114--3125, 2021.

% --- Y ---

\bibitem{xu2024}
Xu, Y., et al. (2024).
\newblock Kinetic analysis of alpha-synuclein aggregation.
\newblock \emph{Nature Communications}, 15, 7083.

\bibitem{yan2023nev}
S.~Yan, X.~Jiang, G.~Tofaris, et al., ``Neuronal extracellular vesicle alpha-synuclein as a biomarker for prodromal Parkinson's disease,'' \textit{JAMA Neurology}, vol.~80, pp.~59--67, 2023. doi:~10.1001/jamaneurol.2022.3924.

\bibitem{yang2020}
W.~Yang \textit{et~al.}, ``Current and projected future economic burden of Parkinson's disease in the U.S.,'' \textit{npj Parkinson's Dis.}, vol.~6, p.~15, 2020.

\bibitem{yano2001ppc}
Y.~Yano, S.~L.~Beal, and L.~B.~Sheiner, ``Evaluating pharmacokinetic/pharmacodynamic models using the posterior predictive check,'' \textit{J.\ Pharmacokinetics Pharmacodynamics}, vol.~28, no.~2, pp.~171--192, 2001, doi:10.1023/A:1011555016423.

\bibitem{yngman2022fullrandom} Yngman, Gunnar and Bjug{\aa}rd Nyberg, Henrik and Nyberg, Joakim and Jonsson, E. Niclas and Karlsson, Mats O. ``An introduction to the full random effects model,'' CPT:PSP, vol. 11, no. 2, pp. 149--160, 2022. doi: 10.1002/psp4.12741

\bibitem{yoon2018}
J.~Yoon, J.~Jordon, and M.~van~der~Schaar, ``GAIN: Missing data imputation using generative adversarial nets,'' in \textit{Proc.\ ICML}, 2018, pp.~5689--5698.

\bibitem{you2020}
J.~You, X.~Ma, D.~Y.~Ding, M.~Kochenderfer, and J.~Leskovec, ``Handling missing data with graph representation learning,'' in \textit{Proc.\ NeurIPS}, 2020.

\bibitem{ypma2022}
R.~J.~F.~Ypma \textit{et~al.}, ``Multi-state modelling of Parkinson's disease progression,'' \textit{J.\ Neurol.}, 2022.

% === 2026-04-09: 4-topic parallel research sweep additions (16 new entries) ===
% Source: parallel general-purpose agents with WebFetch/Consensus verification
% Groups: (T1) identifiability/sloppy models, (T2) CSF ELISA cross-reactivity,
%         (T3) prasinezumab/cinpanemab MoA for S5 counterfactual
% All verified per ~/.claude/skills/research discipline

% --- Topic 1: Practical identifiability / sloppy models ---

\bibitem{zheng2019}
Zheng, Y.Q., et al. (2019).
\newblock Local vulnerability and global connectivity jointly shape neurodegenerative disease propagation.
\newblock \emph{PLoS Biology}, 17(11), e3000357.

\bibitem{zhou2021datscan}
Y.~Zhou \textit{et~al.}, ``Robust Bayesian analysis of early-stage Parkinson's disease progression using DaTscan images,'' \textit{IEEE Transactions on Medical Imaging}, 2021. PMC10258684.

% --- Topic 2: CSF α-syn ELISA cross-reactivity (Block 3 P2 CSF prior) ---

\bibitem{zou2005elastic}
H.~Zou and T.~Hastie, ``Regularization and variable selection via the elastic net,'' \textit{J.\ R.\ Stat.\ Soc.\ B}, vol.~67, no.~2, pp.~301--320, 2005.

\end{thebibliography}
